%% toc: P15 164 164.1
\chapter{Evidential Relevance: Ordinal Equivalence, Direction, and the Standing of Merin's Claims}
\label{sec:p15-w02c-p15-027-relevance}

Chapter~\ref{sec:p15-rationality-4-3} sets out Arthur Merin's account of
evidential relevance, and Chapter~\ref{sec:p15-w02b-p15-024-relevance}
supplies the status of its axioms, the provenance of its measures and the
standing of the works it depends on. This chapter supplies further
parts of the same 2021 chapter \cite[pp.~265--275]{Merin2021Relevance}.
It begins with Merin's reasons for stating relevance as a product
condition and his reading of unconditional relevance as the tautological
case of conditional relevance. The first part after that is his reading of conclusive relevance as what an
unreserved conditional assertion conveys, with the qualification he
attaches to it, together with his counterexample to the transitivity of
positive relevance in its general form and his drug-and-cure example
under the assumptions he states for it. The second is the complete form
of his definition of a relevance measure, including the requirement that
its values fill an interval, with the ranges of the seven measures he
lists proved. The third is his definition of ordinal equivalence together
with its second clause, that ordinal equivalence amounts to an increasing
transformation, proved in both directions and bounded by two
counterexamples, his remark on what ordinally equivalent measures can
still differ in, and his remark that the axiomatic characterization fixes
measures only up to ordinal equivalence. The fourth is his argument for a
directed notation and for the word \emph{relevance}, his claim that one
pair of propositions can be more relevant in one direction than the other
under one measure, less under another and equally under a third, his
law that D and, apart from extreme negativity, L favour the less
probable source, and his remark that Carnap's covariance measure is
amount-symmetric although it violates POP-D. A section on negation then
proves his statement that D, L and Z invert under negation of the target
and not of the source. The fifth and sixth restore his statement of the raven
problem and the full argument of his Linda example, with the empirical
report in that example given the standing its evidence allows; a further
section does the same for his reading of base-rate neglect. The seventh
is his mixture theorem for
disjoint disjunctions in the boundary cases that the existing statement
leaves out. The last lists every claim of the chapter by kind, with the
standing of each.

The inputs are finite probability models, one or more source and target
propositions, and the measures of
Definition~\ref{def:p15-r43-measures}. The conclusions are of three
kinds: formal facts about measures, each proved in the annex; the
content of Merin's non-formal claims, stated in independent words with
his locators; and, for each non-formal claim, what the chapter offers
in its support. Notation follows
Definitions~\ref{def:p15-r43-polarity}, \ref{def:p15-r43-extreme}
and~\ref{def:p15-r43-measures}: \(R_T(A\Yright B)_P\) is the relevance
of source \(A\) to target \(B\) under \(P\) on a measure of type
\(T\), \(\iota_T\) is its neutral value, and
\(\Delta_P(A,B)=P(AB)-P(A)P(B)\).

\section{Relevance as a product condition}
\label{sec:p15-w02c-p15-027-product}

Merin reaches the qualitative relations in two steps
\cite[p.~266]{Merin2021Relevance}. His Definition~2 compares the
conditional probability of \(B\) given \(A\) with the probability of
\(B\): \(A\) is positive, irrelevant or negative to \(B\) under \(P\)
according as \(P(B\mid A)>P(B)\), \(P(B\mid A)=P(B)\) or \(P(B\mid A)<P(B)\),
and the irrelevance case carries a further clause, \emph{or \(P(A)=0\)},
because his Definition~1 leaves \(P(B\mid A)\) undefined when
\(P(A)=0\). Multiplying the three comparisons by \(P(A)\) gives his
Definition~3: \(A\) is positive, irrelevant or negative to \(B\)
according as \(P(AB)>P(A)P(B)\), \(P(AB)=P(A)P(B)\) or
\(P(AB)<P(A)P(B)\), which is the sign of \(\Delta_P(A,B)\) used in
Definition~\ref{def:p15-r43-polarity}.

Merin states why he adopts the second form. Definition~3 takes its
intuitive content from Definition~2, which says what it is for learning
\(A\) to raise or lower the credence in \(B\). In return it has two
advantages: it needs no completion clause, and it displays the symmetry
of the relations. His Theorem~1 is that symmetry: \(A\) is positive,
irrelevant or negative to \(B\) exactly when \(B\) is so to \(A\). He
says the proof is evident from Definition~3, and glosses the theorem
intuitively: if \(A\) speaks for \(B\), then \(B\) speaks for \(A\),
however weakly \cite[p.~266]{Merin2021Relevance}.

On the next page he introduces relevance conditional on a proposition.
His Definition~5 calls \(A\) irrelevant to (independent of) \(C\)
conditionally on \(B\) under \(P\), written \((A\perp C\mid B)_P\), when
\(P(AC\mid B)=P(A\mid B)\,P(C\mid B)\); when this holds both for \(B\) and
for \(\neg B\) he writes \((A\perp C\mid\pm B)_P\). He prefaces it with
the probability theorem \(P(A)=P(A\mid\top)\) for any tautology \(\top\),
from which the irrelevance clause of Definition~3 is the special case of
Definition~5 in which the condition is a tautology
\cite[p.~267]{Merin2021Relevance}. His note~3 uses the same pattern for a
conditional polarity, weakening \((A\perp C\mid\pm B)_P\) to
\(A\) being not negative to \(C\) given \(B\) and given \(\neg B\)
\cite[p.~274, n.~3]{Merin2021Relevance}.

\begin{hbtheorem}{thm:p15-w02c-p15-027-product}{proved=annex:p15-w02c-p15-027-product}{Definitions 2, 3 and 5 compared}
Let \(P\) be a probability and \(A,B,C\) events.
\begin{enumerate}
\item Let \(\bowtie\) be any of \(<,=,>\). If \(P(A)>0\), then
  \[
   P(B\mid A)\bowtie P(B)\quad\text{exactly when}\quad
   P(AB)\bowtie P(A)P(B).
  \]
\item If \(P(A)=0\), then \(P(AB)=P(A)P(B)=0\). So Definition~3
  classifies \(A\) as irrelevant to \(B\) with no further clause, and
  this is the classification that the completion clause of
  Definition~2 stipulates. Definitions~2 and~3 therefore agree for every
  \(P\), \(A\) and \(B\).
\item The condition of Definition~3 is unchanged when \(A\) and \(B\) are
  exchanged, so each of the three relations is symmetric for every
  \(P\), including events of probability zero. Without the completion
  clause the ratio form is not symmetric: on three outcomes with masses
  \(0\), \(1/2\), \(1/2\), let \(A\) be the first outcome and \(B\) the
  second. Then \(P(A\mid B)=0=P(A)\), so \(B\) is irrelevant to \(A\),
  while \(P(B\mid A)\) is undefined, so the ratio comparison does not
  classify \(A\) to \(B\) at all.
\item Symmetry holds for polarity, not for amount. Let \(k\) be uniformly
  distributed on the thirty integers \(15,\ldots,44\), and let \(A\) and
  \(B\) say that \(15\) and \(3\) divide \(k\). Then \(A\) is positive to
  \(B\) and \(B\) to \(A\), with
  \(R_D(A\Yright B)_P=2/3\) and \(R_D(B\Yright A)_P=2/15\).
\item If \(P(K)=1\), in particular if \(K\) is the tautology \(\top\),
  then \(P(X\mid K)=P(X)\) for every event \(X\), and
  \(P(AC\mid K)\bowtie P(A\mid K)P(C\mid K)\) exactly when
  \(P(AC)\bowtie P(A)P(C)\). Hence \((A\perp C\mid K)_P\) holds exactly
  when \(A\) is irrelevant to \(C\) in the sense of Definition~3, and
  the same substitution turns each conditional polarity given \(K\) into
  the corresponding unconditional one.
\item If \(P(B)>0\), then \(P(AC\mid B)\bowtie P(A\mid B)P(C\mid B)\)
  exactly when \(P(ABC)\,P(B)\bowtie P(AB)\,P(BC)\). The product form on
  the right is defined for every \(B\) and, when \(P(B)=0\), holds with
  equality.
\end{enumerate}
\end{hbtheorem}

Items~1 and~2 are Merin's first reason in exact form: the multiplied
condition covers the null case that Definition~2 has to add by hand.
Item~3 is his second reason and his Theorem~1, and its counterexample
shows that the completion clause is what keeps the ratio form
symmetric. Item~4 gives the gloss \emph{however weakly} a numerical
sense: in the divisibility example, which Merin uses for amount symmetry
\cite[p.~272]{Merin2021Relevance}, \(A\) entails \(B\) and speaks for it
strongly by D, while \(B\) speaks for \(A\) with a fifth of that amount.
Merin's qualification \emph{intuitively} marks the gloss as a reading of
the formal symmetry, not a further theorem. Item~5 is his tautology
argument; it needs only that the condition be certain under \(P\), not
that it be a tautology. Item~6 is supplied here and is not in Merin's
text: his Definition~5, as printed, is stated with ratio conditionals
and has no clause for a condition of probability zero, and the handbook's
Definition~\ref{def:p15-r43-extreme} accordingly requires \(P(B)>0\).
Multiplying by \(P(B)^2\) would remove that restriction in the same way
that Definition~3 removes the restriction of Definition~2; Merin does
not take that step, and nothing below depends on it.

\section{Conclusive relevance, transitivity and a stratified population}
\label{sec:p15-w02c-p15-027-extreme}

Extreme positivity, \(A\,\mathrm{expos}_P\,B\), requires
\(P(B)<P(B\mid A)=1\) (Definition~\ref{def:p15-r43-extreme}; Merin's
Definition~4, which he describes as Carnap's 1950 notion modified).
Merin gives it an interpretation in three parts
\cite[p.~266]{Merin2021Relevance}. First, clause~(a) states in a theory
of credal probability what a speaker conveys who asserts an indicative
conditional \emph{if \(A\), then \(B\)} without reservations. He
qualifies this at once: the same conditional can also be asserted with
less certainty. Second, extreme positivity is his explication of what
everyday speech mostly means by saying that \(A\) entails \(B\): \(A\) is
a compelling, and in that sense sufficient, reason for \(B\). Third, the
two halves of the definition carry the two halves of that phrase:
\(P(B\mid A)=1\) makes \(B\) conditionally certain, which is the
sufficiency, and \(P(B)<P(B\mid A)\) is what makes \(A\) a reason. The
Dow-and-prime-number example of Chapter~\ref{sec:p15-rationality-4-3}
shows the second half doing work: there the conditional probability is
one, yet the conditional sounds absurd, because the target was already
certain.

The qualification limits the reading. An unreserved assertion corresponds
to conditional certainty; an assertion made with less certainty, read in
the same way, corresponds to \(P(B)<P(B\mid A)<1\), which is positive
relevance without extreme positivity. The interpretation therefore
covers one kind of conditional assertion, not every use of the
indicative conditional, and the chapter offers it as an explication, not
as a report of how speakers use conditionals. Item~3 of the next theorem
shows that the two halves of the definition are independent.

Merin's counterexample to the transitivity of positive relevance is
stated with three assumptions: \(A\,\mathrm{expos}_P\,B\),
\(C\,\mathrm{expos}_P\,B\) and \(P(AC)=0\). From them he concludes that
\(A\) is positive to \(B\) and \(B\) positive to \(C\), the second by
symmetry, while \(A\) is extremely negative, and hence negative, to \(C\);
he adds that the assumptions are easily satisfied
\cite[p.~267]{Merin2021Relevance}.

\begin{hbtheorem}{thm:p15-w02c-p15-027-transitivity}{proved=annex:p15-w02c-p15-027-transitivity}{The transitivity counterexample in general form}
\begin{enumerate}
\item If \(A\,\mathrm{expos}_P\,B\), \(C\,\mathrm{expos}_P\,B\) and
  \(P(AC)=0\), then \(A\) is positive to \(B\), \(B\) is positive to
  \(C\), \(A\,\mathrm{exneg}_P\,C\), and \(A\) is negative to \(C\).
  The conclusion \(A\,\mathrm{exneg}_P\,C\) uses only \(P(A)>0\),
  \(P(C)>0\) and \(P(AC)=0\).
\item The assumptions of item~1 are jointly satisfiable. If \(A\) and
  \(C\) have positive probability, \(P(AC)=0\), and \(B=A\vee C\) has
  probability below one, all three hold. For instance, on three outcomes
  of mass \(1/3\) each, let \(A\) be the first, \(C\) the second and
  \(B\) the union of the two; then
  \(P(B)=2/3<1=P(B\mid A)=P(B\mid C)\).
\item The two conditions in extreme positivity are independent. On four
  outcomes of mass \(1/4\) each, let \(A\) be the first two outcomes.
  If \(B\) is the whole space, \(P(B\mid A)=1\) but \(P(B)=1\), so \(A\)
  is irrelevant to \(B\). If \(B\) is the first outcome alone,
  \(P(B)=1/4<1/2=P(B\mid A)<1\), so \(A\) is positive but not extremely
  positive to \(B\).
\end{enumerate}
\end{hbtheorem}

Merin's own instance takes \(A\) to be that Trump has won, \(C\) that
Clinton has won and \(B\) that the United States has a new president.
These are his historical illustration, from the 2016 election, not a
claim about any present election. The instance satisfies the assumptions
only in a credal state that leaves \(B\) uncertain and gives each win
positive probability; the three-outcome model of item~2 is such a state,
with the third outcome the one in which neither has won. If \(B\) were
certain, \(A\) would be irrelevant to it by Merin's Theorem~3, and the
example would fail.

Irrelevance within each of two groups does not give irrelevance in the
whole population. Merin illustrates this with a drug-and-cure example
\cite[p.~267]{Merin2021Relevance}. He takes \(B\) and \(\neg B\) to be
complementary properties of the members of a population, being British
and not being British; \(A\) is having been given a drug and \(C\) being
cured of a disease. Two assumptions turn these properties into
propositions with credal probabilities. Each proposition is read as the
statement that \(\iota\) has the property, where \(\iota\) is an
individual sampled at random (Merin's letter, unrelated to the neutral
value \(\iota_T\) of a measure); and the credal probabilities are taken to
reflect the statistical frequencies of the properties. Under these
assumptions, he says, his Theorem~5 allows \(A\) to be positive to \(C\)
in the whole population while irrelevant to \(C\) within each of the
two groups, and he gives a causal story that would produce this: the
drug does nothing for the disease, but both its administration and the
rate of spontaneous recovery are higher among the British. He calls this
an instance of Simpson's paradox and refers to Pearl's chapter of the
same handbook.

\begin{hbtheorem}{thm:p15-w02c-p15-027-simpson}{proved=annex:p15-w02c-p15-027-simpson}{A finite population meeting the stated assumptions}
Take a population of one hundred individuals. Fifty are British; of
these, forty were given the drug, forty were cured, and thirty-two were
both given the drug and cured. Of the fifty others, ten were given the
drug, ten were cured, and two were both. Let \(\iota\) be drawn with
equal probability from the hundred, let \(A\), \(B\) and \(C\) say that
\(\iota\) was given the drug, is British and was cured, and let \(P\)
be the resulting sampling measure, so that the probability of each
proposition is the frequency of the property.
\begin{enumerate}
\item \(P(B)=P(A)=P(C)=1/2\), \(P(AB)=P(BC)=2/5\) and \(P(AC)=17/50\).
\item \((A\perp C\mid\pm B)_P\): among the British,
  \(P(AC\mid B)=16/25=P(A\mid B)P(C\mid B)\); among the others,
  \(P(AC\mid\neg B)=1/25=P(A\mid\neg B)P(C\mid\neg B)\).
\item \(A\) is positive to \(B\) and \(B\) to \(C\), with
  \(\Delta_P(A,B)=\Delta_P(B,C)=3/20\), and \(A\) is positive to \(C\),
  with
  \[
   \Delta_P(A,C)=\frac{9}{100}
     =\frac{\Delta_P(A,B)\,\Delta_P(B,C)}{P(B)P(\neg B)},
  \]
  as Theorem~\ref{thm:p15-r43-screening} requires, and
  \(P(C\mid A)=17/25>1/2=P(C)\).
\end{enumerate}
\end{hbtheorem}

The theorem shows that Merin's two assumptions can be met exactly, by a
finite population and uniform sampling, and that his Theorem~5 then
yields the reported pattern. The counts are supplied here; the
equal-mass model of Chapter~\ref{sec:p15-rationality-4-3} is the same
distribution. The limits are those of the assumptions. Equating credence
with frequency is an assumption about the credal state, which the
example does not argue for; for an infinite population a sampling
distribution would have to be specified. The causal story is one way the
frequencies could arise: the probability model is equally compatible
with other causal structures, and it does not establish that the drug is
ineffective. Conditional independence within the groups is a property of
the frequencies, not a causal finding.

Merin's next section begins with a remark about practice: the most
familiar uses of the symmetry of relevance and of its reversal under
negation (his Theorems~1 and~2) presuppose that propositions are sorted
into two kinds, hypotheses or theories, labelled \(H\), and observational
data reports or evidence statements, labelled \(D\) or \(E\)
\cite[p.~267]{Merin2021Relevance}. The theorems themselves use no such
sorting, and Merin argues on the next page that the sorting often
depends on the application (Section~\ref{sec:p15-w02c-p15-027-direction});
the remark records how the theorems are usually applied, as in the
prediction and old-evidence analysis of
Section~\ref{sec:p15-r43-old-evidence}.

\section{Relevance measures and the interval requirement}
\label{sec:p15-w02c-p15-027-measures}

For propositions \(A,B\) in the domain of \(P\), write
\[
 \pi_P(A,B)=\bigl(P(AB),\,P(A\neg B),\,P(\neg AB),\,P(\neg A\neg B)\bigr)
\]
for their \emph{joint table}, a point of the set \(\Sigma\) of
nonnegative quadruples summing to one. The order of the cells matters:
\(\pi_P(B,A)\) exchanges the two middle cells.

\begin{hbdefinition}{def:p15-w02c-p15-027-measure}{Probabilistic relevance measure}
Merin's Definition~8 calls a function \(R_T(A\Yright B)_P\), defined
for every probability measure \(P\) and every pair \(A,B\) of
propositions in its domain and taking values in the extended reals
\(\mathbb R^*=\mathbb R\cup\{-\infty,\infty\}\), a
\emph{probabilistic relevance measure} when it satisfies two conditions
\cite[p.~269]{Merin2021Relevance}.
\begin{enumerate}
\item[(a)] \emph{Transparency.} The value is a differentiable function
  of the distribution of \(P\) over the four atoms \(AB\),
  \(A\wedge\neg B\), \(\neg A\wedge B\), \(\neg A\wedge\neg B\):
  there is a function \(\phi_T:\Sigma\to\mathbb R^*\) with
  \(R_T(A\Yright B)_P=\phi_T(\pi_P(A,B))\) for all \(P,A,B\).
  In addition, the set of values the function takes is an interval of
  \(\mathbb R^*\), proper or improper.
\item[(b)] \emph{Quality Representation.} There is a value
  \(\iota_T\) in that range such that \(R_T(A\Yright B)_P\) is greater
  than, equal to or less than \(\iota_T\) exactly when \(A\) is
  respectively positive, irrelevant or negative to \(B\) under \(P\).
\end{enumerate}
Merin's note~6 explains the label: transparency means that the measure
holds no unavoidable surprises, and hidden propositions or jumps would
be surprises \cite[p.~274, n.~6]{Merin2021Relevance}.
\end{hbdefinition}

The text of Chapter~\ref{sec:p15-rationality-4-3} states the first half
of condition~(a) and condition~(b). The second half of~(a), the interval
requirement, is a separate restriction on the set of attained values,
and it is not implied by the rest: a function of the table can satisfy
Quality Representation and still take only isolated values. The
differentiability half of~(a) is an interior requirement; as
Theorem~\ref{thm:p15-r43-measures} shows, Z is not differentiable at
its neutral point unless the prior is one half, so the requirement
cannot be read as global differentiability of every listed measure.

\begin{hbtheorem}{thm:p15-w02c-p15-027-ranges}{proved=annex:p15-w02c-p15-027-ranges}{The seven listed measures satisfy the interval requirement}
With the completions of Definition~\ref{def:p15-r43-measures}, the sets
of values attained over all finite probability models and all pairs of
propositions are exactly
\[
\begin{array}{c|ccccccc}
T & D & Q & L & QL & LL & KO & Z\\\hline
\text{range} & (-1,1) & [0,\infty) & [0,\infty] & [-\infty,\infty)
  & [-\infty,\infty] & [-1,1] & [-1,1]
\end{array}
\]
Each is an interval of \(\mathbb R^*\). Every value in each interval is
already attained by a pair with \(0<P(A),P(B)<1\) on a four-point
sample space.

By contrast, the polarity indicator \(\operatorname{sgn}R_D\), with
values \(1,0,-1\) for positive, irrelevant and negative, is a function
of the joint table and satisfies Quality Representation with
\(\iota=0\), but its range \(\{-1,0,1\}\) is not an interval. It is
therefore not a relevance measure in the sense of
Definition~\ref{def:p15-w02c-p15-027-measure}.
\end{hbtheorem}

These ranges agree with those Merin states for the measures
\cite[p.~270]{Merin2021Relevance}. The point of the interval
requirement becomes visible in the next section: when two relevance
measures are increasing transforms of each other, the requirement forces
the transformation to be continuous.

\section{Ordinal equivalence and increasing transformations}
\label{sec:p15-w02c-p15-027-ordinal}

\begin{hbdefinition}{def:p15-w02c-p15-027-ordinal}{Ordinal equivalence}
Merin's Definition~9 calls measure types \(T\) and \(T'\)
\emph{ordinally equivalent}, written \(T\sim_{\uparrow}T'\), when, for
every probability measure \(P\) and all propositions \(A,B,C,D\) in its
domain,
\[
 R_T(A\Yright B)_P\bowtie R_T(C\Yright D)_P
 \quad\text{exactly when}\quad
 R_{T'}(A\Yright B)_P\bowtie R_{T'}(C\Yright D)_P,
\]
where \(\bowtie\) is any one of \(<,=,>\). He adds that, equivalently,
\(T\sim_{\uparrow}T'\) holds exactly when some strictly increasing
(\emph{isotonic}) function \(f\) gives \(R_{T'}=f(R_T)\); his note~9
calls a function \(f:\mathbb R\to\mathbb R\) isotonic when \(a<b\)
implies \(f(a)<f(b)\) \cite[p.~270 and p.~274, n.~9]{Merin2021Relevance}.
\end{hbdefinition}

Two features of the definition matter. Its comparisons are made within
one probability measure \(P\), so the definition does not by itself
compare a value attained under one model with a value attained under
another. And the function in note~9 is defined on the real line, while
several of the measures take infinite values. The theorem below shows
that the first feature is harmless for transparent measures and that the
second requires the function to be taken on the set of attained values.

\begin{hbtheorem}{thm:p15-w02c-p15-027-ordinal}{proved=annex:p15-w02c-p15-027-ordinal}{Ordinal equivalence is increasing transformation on the range}
Let \(T,T'\) be measure types, and write \(\operatorname{ran}T\) for
the set of values \(R_T\) attains.
\begin{enumerate}
\item If \(f\) is strictly increasing on \(\operatorname{ran}T\) and
  \(R_{T'}(A\Yright B)_P=f(R_T(A\Yright B)_P)\) for all \(P,A,B\), then
  \(T\sim_{\uparrow}T'\).
\item If \(T\) and \(T'\) are both functions of the joint table, as in
  the first half of Definition~\ref{def:p15-w02c-p15-027-measure}(a),
  and \(T\sim_{\uparrow}T'\), then there is exactly one function
  \(f:\operatorname{ran}T\to\mathbb R^*\) with
  \(R_{T'}=f(R_T)\) everywhere. It is strictly increasing and maps
  \(\operatorname{ran}T\) onto \(\operatorname{ran}T'\).
\item If moreover both ranges are intervals, as the second half of
  Definition~\ref{def:p15-w02c-p15-027-measure}(a) requires, this \(f\)
  is continuous on \(\operatorname{ran}T\) in the order topology of
  \(\mathbb R^*\).
\item Without the table condition, item~2 fails. Let \(c(P)=1\) when
  the sample space of \(P\) has at most four points and \(c(P)=2\)
  otherwise, and put \(R_{T'}(A\Yright B)_P=c(P)\,R_D(A\Yright B)_P\).
  Then \(T'\) satisfies Quality Representation with neutral value zero,
  its range is an interval, and \(D\sim_{\uparrow}T'\) under
  Definition~\ref{def:p15-w02c-p15-027-ordinal}, but no function \(f\)
  satisfies \(R_{T'}=f(R_D)\).
\item The function of item~2 need not extend to a strictly increasing
  function on \(\mathbb R\) or on \(\mathbb R^*\). L and LL are
  ordinally equivalent, with \(R_{LL}=\ln R_L\) and
  \(\ln(+\infty)=+\infty\), and L attains \(+\infty\), so no
  \(f:\mathbb R\to\mathbb R\) gives \(R_{LL}=f(R_L)\). Nor does any
  strictly increasing \(f:\mathbb R^*\to\mathbb R^*\), since it would
  need \(f(0)=-\infty\) and a value below \(-\infty\) at \(-1\).
\end{enumerate}
\end{hbtheorem}

So Merin's second clause holds, for transparent measures, when \(f\) is
read as a function on the attained values. That reading is the one his
own examples need: KO, QL and LL are presented as transforms of L and
Q \cite[pp.~269--270 and p.~274, n.~7]{Merin2021Relevance}, and each
transform is defined at the infinite values the measure takes. The
transparency requirement is what lets one function serve all models;
item~4 shows a measure that is ordinally equivalent to D in every single
model but rescales differently in different models. Item~3 is the
reason the interval requirement is not idle: two relevance measures that
are ordinally equivalent differ by a continuous increasing relabelling of
values, as KO and LL do through \(R_{KO}=\tanh(\tfrac12R_{LL})\).

Merin adds, right after Definition~9, that measures which are distinct
but ordinally equivalent can still differ in two ways: in how convenient
their ranges are, and in their algebraic properties, the second being in
his view the more important for applications in cognitive engineering
\cite[p.~270]{Merin2021Relevance}. Both differences are visible within
the class of L. On ranges, L takes values in \([0,\infty]\) with neutral
value one, LL in \([-\infty,\infty]\) with neutral value zero, and KO in
the bounded interval \([-1,1]\)
(Theorem~\ref{thm:p15-w02c-p15-027-ranges}). On algebra, conditionally
independent sources compose by multiplication under L and by addition
under LL, while KO composes by the addition law of the hyperbolic tangent,
and among increasing transforms of L only positive powers keep the
product rule and only positive multiples of LL keep the sum rule
(Theorem~\ref{thm:p15-w02b-p15-024-ko-transform}); negating the target
inverts L and negates LL (Theorem~\ref{thm:p15-r43-measures}), a
reciprocal under L that is a sign change under LL. These are formal
facts. The weighting, that algebraic properties matter more than ranges,
is Merin's judgment; the chapter names no cognitive-engineering
application and gives no argument for it beyond the conditional
additivity of LL that he cites in recommending it as the default
(p.~274).

The same limit applies to the axiomatic characterization. Merin states
it, his Theorem~7, after Crupi, Chater and Tentori (2013) and Crupi and
Tentori (2016), for a real-valued function that satisfies Transparency
with the possible exception of its smoothness requirement: such a
function is D, Q or L, or an isotonic transformation of one of them,
exactly when it satisfies POP-D together with MDHD, LIK-D or MCIE
respectively (Theorem~\ref{thm:p15-r43-characterization}). He then says
what the theorem does not do. Its \emph{if} parts yield ordinal
equivalence classes \([T]=\{T'\mid T'\sim_{\uparrow}T\}\), and to
determine the specific measures D, Q, L, QL, LL, KO and Z further
conditions are needed, smoothness among them
\cite[p.~271]{Merin2021Relevance}. The chapter names no condition other
than smoothness. The next theorem shows, within the class of D, what
smoothness excludes and what it leaves.

\begin{hbtheorem}{thm:p15-w02c-p15-027-classes}{proved=annex:p15-w02c-p15-027-classes}{The axioms fix a class, not a measure}
Compare measures on the comparison domain of
Theorem~\ref{thm:p15-r43-characterization}, where source and target
probabilities lie strictly between zero and one.
\begin{enumerate}
\item If \(f\) is strictly increasing on \((-1,1)\), then
  \(f\circ R_D\) satisfies POP-D and MDHD, and \(f\circ R_D\sim_{\uparrow}D\).
\item Let \(R_{D^3}=(R_D)^3\). It is differentiable as a function of the
  joint table wherever all four cells are positive, its range is
  \((-1,1)\), its neutral value is zero, and it satisfies POP-D and MDHD;
  yet it is not D: where \(P(H)=1/4\) and \(P(H\mid E)=3/4\),
  \(R_D(E\Yright H)_P=1/2\) and \(R_{D^3}(E\Yright H)_P=1/8\).
\item Let \(R_{D^\sharp}=R_D\) where \(R_D\geq0\) and
  \(R_{D^\sharp}=2R_D\) where \(R_D<0\). It satisfies POP-D and MDHD,
  satisfies Quality Representation with neutral value zero, and has the
  interval \((-2,1)\) as range. At the joint table with all four cells
  equal to \(1/4\), where \(A\) is irrelevant to \(B\), it is not
  differentiable: moving mass \(t\) from the cell \(A\neg B\) to the cell
  \(AB\) changes it at rate one for \(t>0\) and at rate two for \(t<0\).
\end{enumerate}
\end{hbtheorem}

Item~1 is the reason the axioms cannot single out D: they are
statements about order, and an increasing relabelling keeps every order
statement. Item~3 shows that smoothness does real work, since it removes
from the class a member that the axioms, the interval requirement and
Quality Representation all allow. Item~2 shows that smoothness is not
enough: \(R_{D^3}\) shares with D the axioms, smoothness, the range and
the neutral value, and still differs from it. So Merin's \emph{further
conditions} must include more than smoothness, the range and the neutral
value; which conditions he intends the chapter does not say, and none is
supplied here. Within the class of L the three measures L, LL and KO
already illustrate the point, with three different ranges
(Theorem~\ref{thm:p15-w02c-p15-027-ranges}).

\section{The direction of measurement can be ranked three ways}
\label{sec:p15-w02c-p15-027-direction}

Merin motivates his signpost notation with a claim about direction. For
two mutually relevant propositions, the relevance of the first to the
second can exceed the relevance of the second to the first on one
measure, fall short of it on another and match it on a third
\cite[p.~268]{Merin2021Relevance}.
Theorem~\ref{thm:p15-r43-measures} gives the comparison for D and the
symmetry of Q. The following theorem supplies the general formulas and
a single model in which all three outcomes occur. The third measure used
there, \(\delta p(A\Yright B)=P(B\mid A)-P(B\mid\neg A)\), is the
quantity Merin reports as an explication of participants' relevance
assessments \cite[p.~273]{Merin2021Relevance}; it is used only on
\(0<P(A)<1\), where it is defined, and Merin gives it no completion.

\begin{hbtheorem}{thm:p15-w02c-p15-027-direction}{proved=annex:p15-w02c-p15-027-direction}{Direction comparisons under D, L, Z, Q and \texorpdfstring{\(\delta p\)}{delta p}}
Let \(0<P(A),P(B)<1\) and \(\delta=\Delta_P(A,B)\ne0\). Then
\[
\begin{aligned}
 R_D(A\Yright B)_P&=\frac{\delta}{P(A)},\qquad
 \delta p(A\Yright B)=\frac{\delta}{P(A)(1-P(A))},\\
 R_Q(A\Yright B)_P&=1+\frac{\delta}{P(A)P(B)},\\
 R_Z(A\Yright B)_P&=
 \begin{cases}
  \dfrac{\delta}{P(A)(1-P(B))},&\delta>0,\\[1ex]
  \dfrac{\delta}{P(A)P(B)},&\delta<0.
 \end{cases}
\end{aligned}
\]
Consequently Q, and Z on negative relevance, give the two directions
equal values; D, and Z on positive relevance, give the larger absolute
value to the direction whose source is less probable; and \(\delta p\)
gives the larger absolute value to the direction whose source has the
smaller \(P(1-P)\), that is, the source whose probability is further
from one half.

In the model with cell masses
\[
 P(AB)=\tfrac{17}{50},\quad P(A\neg B)=\tfrac{3}{50},\quad
 P(\neg AB)=\tfrac{28}{50},\quad P(\neg A\neg B)=\tfrac{2}{50},
\]
\(A\) is negative to \(B\), \(P(A)=2/5\), \(P(B)=9/10\), and
\[
\begin{array}{c|ccccc}
 & D & L & \delta p & Q & Z\\\hline
 A\Yright B & -1/20 & 17/27 & -1/12 & 17/18 & -1/18\\
 B\Yright A & -1/45 & 51/56 & -2/9 & 17/18 & -1/18
\end{array}
\]
Measured by distance from the neutral value, the relevance of \(A\) to
\(B\) exceeds that of \(B\) to \(A\) under D and under L, falls short
of it under \(\delta p\), and equals it under Q and under Z.
\end{hbtheorem}

The model realizes Merin's claim with measures the chapter itself uses.
It also shows why his notation records direction rather than content
type: nothing about \(A\) or \(B\) marks either as evidence or
hypothesis, yet the ranking of the two directions depends on the
measure. For L the comparison is made on the logarithmic scale, where
\(|\ln(17/27)|>|\ln(51/56)|\). The symmetry of Z on negative
relevance is the partial-inconsistency axiom of
Section~\ref{sec:p15-w02b-p15-024-dependents}.

Merin states the D comparison as a law: if \(A\) and \(B\) are not
independent under \(P\), the absolute value of
\(R_D(A\Yright B)_P\) is greater than, equal to or less than that of
\(R_D(B\Yright A)_P\) according as \(P(A)\) is less than, equal to or
greater than \(P(B)\); and he adds that L behaves in the same way except
for extreme negativity \cite[p.~272]{Merin2021Relevance}. The display on
that page prints the same three-way comparison sign on both sides of
the biconditional; read with the two signs matched part for part it
would give the larger D value to the more probable source, which
contradicts the sentence it formalizes and the formula of
Theorem~\ref{thm:p15-w02c-p15-027-direction}. The statement above follows
the sentence. For L, distance from the neutral value one is measured by
\(|R_{LL}|=|\ln R_L|\), with \(|\ln 0|=|\ln(+\infty)|=+\infty\). The next
theorem computes both L values in every case, including the conclusive
ones, and locates the exception.

\begin{hbtheorem}{thm:p15-w02c-p15-027-likewise}{proved=annex:p15-w02c-p15-027-likewise}{L compares the two directions as D does, except at extreme negativity}
Let \(A\) and \(B\) not be independent under \(P\), and write
\(u=P(A)\), \(v=P(B)\), \(x=P(AB)\). Then \(0<u,v<1\),
\(x\ne uv\) and \(\max\{0,u+v-1\}\leq x\leq\min\{u,v\}\), and exactly
one of the following holds.
\begin{enumerate}
\item \emph{Interior.} If \(0<x<\min\{u,v\}\), then
  \[
   R_L(A\Yright B)_P=\frac{x(1-v)}{v(u-x)},\qquad
   R_L(B\Yright A)_P=\frac{x(1-u)}{u(v-x)},
  \]
  both finite and positive, and
  \[
   R_L(A\Yright B)_P-R_L(B\Yright A)_P=\frac{x(v-u)(x-uv)}{uv(u-x)(v-x)}.
  \]
\item \emph{Extreme positivity in one direction.} If \(x=u<v\), then
  \(A\,\mathrm{expos}_P\,B\), and
  \[
   R_L(A\Yright B)_P=+\infty,\qquad
   R_L(B\Yright A)_P=\frac{1-u}{v-u}\in(1,+\infty).
  \]
  In the mirror case \(x=v<u\), \(R_L(B\Yright A)_P=+\infty\) and
  \(R_L(A\Yright B)_P=(1-v)/(u-v)\in(1,+\infty)\).
\item \emph{Extreme positivity in both directions.} If \(x=u=v\), then
  \(R_L(A\Yright B)_P=R_L(B\Yright A)_P=+\infty\).
\item \emph{Extreme negativity.} If \(x=0\), then \(A\,\mathrm{exneg}_P\,B\)
  and \(B\,\mathrm{exneg}_P\,A\), and
  \(R_L(A\Yright B)_P=R_L(B\Yright A)_P=0\).
\end{enumerate}
In cases~1 to~3,
\(|\ln R_L(A\Yright B)_P|\) is greater than, equal to or less than
\(|\ln R_L(B\Yright A)_P|\) according as \(u\) is less than, equal to or
greater than \(v\). In case~4 the two distances are both \(+\infty\)
whatever \(u\) and \(v\) are, while
\(|R_D(A\Yright B)_P|=v\) and \(|R_D(B\Yright A)_P|=u\).
\end{hbtheorem}

So Merin's exception is exact: the D pattern carries over to L in every
case but extreme negativity, including both conclusive positive cases,
where the source that entails the other up to probability zero has
infinite L and the reverse direction a finite value above one. For
instance, with \(u=1/4\), \(v=1/2\) and \(x=1/4\), as when \(A\) entails
\(B\), \(R_L(A\Yright B)_P=+\infty\) and
\(R_L(B\Yright A)_P=(3/4)/(1/4)=3\); with \(x=0\) instead, both values
are zero, while D still separates the directions, at \(-1/2\) and
\(-1/4\). Merin's divisibility example on the same page is of the first
kind: A, that fifteen divides \(k\), entails B, that three divides
\(k\).

Amount symmetry, Merin's Definition~7, requires
\(R_T(B\Yright A)_P=R_T(A\Yright B)_P\) for all \(A\), \(B\) and \(P\)
\cite[p.~269]{Merin2021Relevance}. Among the measures he discusses, Q and
QL have it. In his features-in-use subsection he adds a third:
Carnap's measure \(R_C(E\Yright H)_P=P(EH)-P(E)P(H)\), which is
amount-symmetric but violates POP-D \cite[p.~272]{Merin2021Relevance}.
Theorem~\ref{thm:p15-w02b-p15-024-carnap-bridge} proves the violation of
POP-D. The next theorem proves the amount symmetry, shows in one model
both the failure of POP-D and the failure of LIK-D, and contrasts \(R_C\)
with Q, the amount-symmetric measure that satisfies POP-D.

\begin{hbtheorem}{thm:p15-w02c-p15-027-covariance}{proved=annex:p15-w02c-p15-027-covariance}{Carnap's measure is amount-symmetric}
\begin{enumerate}
\item For every \(P\), \(E\) and \(H\),
  \[
   R_C(E\Yright H)_P=\Delta_P(E,H)=R_C(H\Yright E)_P .
  \]
  In terms of the joint table,
  \[
   R_C=P(EH)\,P(\neg E\neg H)-P(E\neg H)\,P(\neg EH).
  \]
  It is a polynomial in
  the cells, its range is the interval \([-1/4,1/4]\), and it satisfies
  Quality Representation with neutral value zero; so it is a relevance
  measure in the sense of Definition~\ref{def:p15-w02c-p15-027-measure}.
\item On forty equally likely outcomes, let \(E\) be outcomes
  \(1\)--\(20\), \(H_1\) outcomes \(1\)--\(4\), and \(H_2\) outcomes
  \(5\)--\(19\) and \(21\)--\(25\). Then \(P(E)=P(H_2)=1/2\),
  \(P(H_1)=1/10\), \(P(E\mid H_1)=1\) and \(P(E\mid H_2)=3/4\), while
  \[
  \begin{aligned}
   R_C(E\Yright H_1)_P&=R_C(H_1\Yright E)_P=\tfrac1{20},\\
   R_C(E\Yright H_2)_P&=R_C(H_2\Yright E)_P=\tfrac18 .
  \end{aligned}
  \]
  With \(E\) as source, the larger likelihood gets the smaller value, so
  LIK-D fails. With \(E\) as target, the source giving \(E\) the larger
  posterior probability gets the smaller value, so POP-D fails.
\item In the same model, \(R_Q(H_1\Yright E)_P=2>3/2=R_Q(H_2\Yright E)_P\).
  So the two amount-symmetric measures Q and \(R_C\) order the same two
  sources oppositely for the same target, and they are not ordinally
  equivalent.
\end{enumerate}
\end{hbtheorem}

Item~2 shows how amount symmetry ties the two axioms together: by
symmetry, a failure of LIK-D for a source is a failure of POP-D for the
same pair read the other way. This is the contrapositive face of the
first item of Theorem~\ref{thm:p15-w02b-p15-024-axiom-status}, which
derives LIK-D from POP-D and amount symmetry; \(R_C\) shows that amount
symmetry contributes nothing in that derivation without POP-D. Since
\(R_C(E\Yright H)_P=P(E)\,R_D(E\Yright H)_P\), it weights D by the
probability of the source; in the model the more probable source \(H_2\)
receives the larger value although it raises the probability of \(E\)
less. Merin draws no further conclusion from the example beyond
listing \(R_C\) as amount-symmetric and non-POP-D. The composition laws
of \(R_C\) that he criticizes are proved in
Theorem~\ref{thm:p15-r43-composition}, and his evaluation of them is item~7
of the normative list in Section~\ref{sec:p15-w02c-p15-027-claims}.

Merin draws two conclusions from the direction claim, one about notation
and one about words, in his subsection on notation and terminology
\cite[p.~268]{Merin2021Relevance}. The argument has four steps.
\begin{enumerate}
\item The literature commonly glosses a degree of relevance as an amount
  of incremental confirmation of a hypothesis \(H\) by evidence \(E\)
  under \(P\) and a measure type, written \(C_P(H,E)\) and sometimes
  \(C_P(E,H)\). The comma in this format gives no natural cue to
  direction.
\item The gloss does suggest a direction, from evidence to hypothesis.
  But one often reasons from hypotheses to possible data, and often
  there is no warrant for classifying the contents as evidence and
  hypothesis at all. Relying on the letters \(E\) and \(H\) therefore
  invites one to identify a proposition's position in the argument with
  the type of its content too early.
\item Along with this, the words \emph{confirmation},
  \emph{corroboration} and \emph{support} face two difficulties. Without
  the prefix \emph{incremental} they are easily misread as naming the
  conditional probability \(P(H\mid E)\); and with the prefix
  \emph{negative} they sound awkward. The word \emph{relevance}, used by
  Keynes (1921), Carnap (1950) and Jeffrey (2004), has neither
  difficulty.
\item So a notation that is neutral about content is preferable, at least
  beside the usual philosophy-of-science reading; Merin proposes the
  signpost separator, which marks the conventional direction of
  measurement as the words \emph{relevance of \(A\) to \(B\)} do, and
  calls the two arguments source and target (his Definition~6).
\end{enumerate}
The model of Theorem~\ref{thm:p15-w02c-p15-027-direction} gives the first
difficulty of step~3 a numerical form: there
\(P(B\mid A)=17/20\), so \(B\) is highly probable given \(A\), yet \(A\)
is negative to \(B\), since \(P(B)=9/10\). Saying that \(A\) gives \(B\)
strong support would be true in the conditional-probability sense and
false in the incremental one. Steps~1 to~3 are claims about usage and
about how words sound, made without data; step~4 is a recommendation
that rests on them and on the direction claim, which is the only part
given a formal basis. Section~\ref{sec:p15-w02b-p15-025-sources}
records the first difficulty of step~3 in its comparison of what
different sources measure.

Two further remarks frame the proposal. Before the argument, Merin says
that a helpful notation should show the intended direction at once,
without tacitly relying on a classification of contents as evidence or
hypothesis, which often depends on the application; and that, with such
a notation, what motivates the conventional assignments of direction can
then be investigated freely \cite[p.~268]{Merin2021Relevance}. This is a
methodological claim about what the notation makes possible. The chapter
does not carry out such an investigation in general. It does fix
directions case by case, and one case shows the separation at work: on
the representativeness reading of the Linda task, the description
\(E\), which its letter marks as evidence, stands in target position
(Section~\ref{sec:p15-w02c-p15-027-linda}). After his Definition~8 Merin
adds that he will mostly use the format \(R(E\Yright H)\), which
identifies the source with evidence \(E\) and the target with a
hypothesis \(H\), because it eases cross-reference to the confirmation
literature, as the letters \(A\) and \(B\) do for the content-neutral
mathematical literature \cite[p.~269]{Merin2021Relevance}. This is a
choice of presentation, not a retraction of step~2: in the signpost
format the letters name argument positions, and the position, not the
letter, fixes the direction. This book follows the same practice.

\section{Negation in target and in source position}
\label{sec:p15-w02c-p15-027-negation}

The first homomorphism Merin discusses is relevance-functionality of
negation: negating a proposition sends the relevance value \(r\) to
\(-r\) when the neutral value is zero and to \(1/r\) when it is one.
Target inversion says that \(E\) always speaks for \(H\) exactly as
strongly as it speaks against \(\neg H\); source inversion is the corresponding
law for \(E\) and \(\neg E\) with \(H\) fixed. Merin reports that professed
opinion usually commends the first unconditionally and recommends
denying the second, and he states that D, L and Z share both the
positive and the negative property \cite[p.~272]{Merin2021Relevance}.
Theorem~\ref{thm:p15-r43-measures} states target inversion for the
three measures and says that source negation need not act in the same
way; its proof gives the D formula and an instance for L. The theorem below proves the pair for all three, including the
completions, characterizes exactly when source inversion holds for
each, and gives the instance for Z. Write \(\mathrm{inv}_T(r)=-r\) for
\(T=D,Z\) and \(\mathrm{inv}_L(r)=1/r\), with \(1/0=+\infty\) and
\(1/(+\infty)=0\).

\begin{hbtheorem}{thm:p15-w02c-p15-027-negation}{proved=annex:p15-w02c-p15-027-negation}{D, L and Z invert under target negation and not under source negation}
Let \(T\in\{D,L,Z\}\), with the completions of
Definition~\ref{def:p15-r43-measures}.
\begin{enumerate}
\item \emph{Target.} For every \(P\), \(E\) and \(H\),
  \(R_T(E\Yright\neg H)_P=\mathrm{inv}_T\bigl(R_T(E\Yright H)_P\bigr)\).
\item \emph{Degenerate source or target.} If \(P(E)\in\{0,1\}\) or
  \(P(H)\in\{0,1\}\), then \(R_T(E\Yright H)_P\) and
  \(R_T(\neg E\Yright H)_P\) both equal the neutral value, so
  \(R_T(\neg E\Yright H)_P=\mathrm{inv}_T\bigl(R_T(E\Yright H)_P\bigr)\).
\item \emph{Source.} Let \(0<e<1\) and \(0<h<1\), where \(e=P(E)\),
  \(h=P(H)\), \(q=P(H\mid E)\), \(a=P(E\mid H)\), \(b=P(E\mid\neg H)\).
  Then \(R_T(\neg E\Yright H)_P=\mathrm{inv}_T\bigl(R_T(E\Yright H)_P\bigr)\)
  holds
  \begin{itemize}
  \item for D if and only if \(q=h\) or \(e=1/2\), because
    \(R_D(\neg E\Yright H)_P=-\dfrac{e}{1-e}\,R_D(E\Yright H)_P\);
  \item for L if and only if \(a=b\) or \(a+b=1\);
  \item for Z if and only if \(q=h\), or \(q>h\) and \(e=h\), or
    \(q<h\) and \(e+h=1\).
  \end{itemize}
\item \emph{Instances.} Take the cell masses
  \[
   P(EH)=\tfrac{3}{16},\quad P(E\neg H)=\tfrac{1}{16},\quad
   P(\neg EH)=\tfrac{5}{16},\quad P(\neg E\neg H)=\tfrac{7}{16},
  \]
  so that \(P(H)=1/2\), \(P(E)=1/4\), \(P(H\mid E)=3/4\) and
  \(P(H\mid\neg E)=5/12\). Then
  \[
  \begin{aligned}
   R_Z(E\Yright H)_P&=\tfrac12,\qquad
   R_Z(E\Yright\neg H)_P=-\tfrac12,\\
   R_Z(\neg E\Yright H)_P&=-\tfrac16\neq-\tfrac12=-R_Z(E\Yright H)_P,
  \end{aligned}
  \]
  and likewise \(R_D(\neg E\Yright H)_P=-1/12\ne-1/4\) and
  \(R_L(\neg E\Yright H)_P=5/7\ne1/3\). With cell masses
  \(3/16\), \(5/16\), \(1/16\), \(7/16\) in the same order, so that
  \(P(E)=1/2\) and \(P(H)=1/4\), D inverts under source negation
  (\(1/8\) and \(-1/8\)) but Z does not (\(1/6\) and \(-1/2\)), nor
  does L (\(9/5\) and \(3/7\)). With cell masses \(1/8\), \(1/8\),
  \(1/8\), \(5/8\), so that \(P(E)=P(H)=1/4\), Z inverts (\(1/3\) and
  \(-1/3\)) but D does not (\(1/4\) and \(-1/12\)), nor does L
  (\(3\) and \(3/5\)).
\end{enumerate}
\end{hbtheorem}

The theorem bears out Merin's statement as a statement about all
probability functions: target inversion holds without condition for the
three measures, and source inversion fails for each in some model.
The conditions under which source inversion does hold differ between
the three, so the shared negative property is a failure of the
unconditional law, not a common set of exceptions. For Z with
\(P(H)=1/2\) the two branches coincide and \(R_Z=2R_D\), so the first
instance fails for Z for the same reason as for D, the factor
\(e/(1-e)=1/3\); the second instance shows that Z can fail even where
D inverts, because the source and its negation fall on different
branches, which divide by \(1-h\) and by \(h\). Merin cites Eells and
Fitelson (2002) for the professed opinion and gives no argument of his
own for preferring inversion in target position and its denial in
source position; the theorem establishes which measures have the pair,
not that the pair is desirable.

\section{The raven problem as Merin states it}
\label{sec:p15-w02c-p15-027-raven}

Merin sets out the raven problem in three steps
\cite[pp.~271--272]{Merin2021Relevance}. The universal hypothesis \(H\),
that every raven is black, receives support from the conjunction
\(A=A_1A_2\) about a randomly encountered entity \(a\), where \(A_1\)
says that \(a\) is a raven and \(A_2\) that \(a\) is black. Its logical
equivalent \(G\), that all nonblack things are nonravens, is supported
by \(B=B_1B_2\), where \(B_1\) says that \(a\) is not black and \(B_2\)
that \(a\) is not a raven. In the third step, \(C\) says that \(a\) is a
white shoe. As printed, this step says that \(G\) is entailed by \(C\)
and thereby supported. Given that white things are not black and shoes
are not ravens, what \(C\) entails is the specimen description \(B\),
not the universal \(G\); note~15 calls the third step an inessential
embellishment that makes things vivid
\cite[p.~272 and p.~274, n.~15]{Merin2021Relevance}.

The problem, in Merin's statement, is this. A deductivist account of
support that seems reasonable requires evidence to bear in the same way
on hypotheses that are logically equivalent. So \(B\), by supporting \(G\), should support \(H\), and with
mild further assumptions \(C\) should as well; and that verdict is
counterintuitive. Merin does not state the further assumptions for
\(C\). The probabilistic proposals he reports accept the transfer
between equivalents and locate the difference elsewhere: nonblack
things vastly outnumber ravens, so \(B\) is far more probable and less
surprising than \(A\), and with some irrelevance assumptions, such as
Hosiasson-Lindenbaum's independence of \(H\) from each of \(A_1\) and
\(B_1\), both \(A\) and \(B\) come out positive to \(H\), \(B\) only
negligibly. Merin then claims that this negligible increment accounts
for the intuition, held before any theory, that \(B\) has no bearing on
\(H\), and that the account carries over to \(C\).

\begin{hbtheorem}{thm:p15-w02c-p15-027-equivalents}{proved=annex:p15-w02c-p15-027-equivalents}{Probabilistic relevance satisfies the equivalence condition}
If \(G\) and \(H\) are logically equivalent, or more generally if
\(P(G\mathbin{\triangle}H)=0\), then for every measure type \(T\) that is
a function of the joint table and every source \(B\),
\(R_T(B\Yright G)_P=R_T(B\Yright H)_P\). In the explicit raven model of
Theorem~\ref{thm:p15-r43-rankings}, \(R_L(A\Yright H)_P=2\) and
\(R_L(B\Yright H)_P=160/159\), so \(B\) is positive to \(H\) and to
\(G\) by the same amount, and that amount is \(1/159\) above the
neutral value one.
\end{hbtheorem}

The theorem fixes the shape of the probabilistic answer. It cannot deny
that \(B\) supports \(H\) once it grants that \(B\) supports \(G\),
because it assigns the same value to both. Its answer to the problem
is therefore not that the deductivist transfer fails but that the
transferred support is tiny, and Theorem~\ref{thm:p15-r43-rankings}
makes the ordering of \(A\) above \(B\) the same under every measure
that satisfies POP-D. Merin's further claim, that the tiny increment
explains the intuition of irrelevance, is an explanatory hypothesis
about intuitive judgment: it presupposes that people register a
negligible positive increment as no support. The chapter gives no data
on such judgments, and the formal results do not supply any. For the
white shoe, the relevance of \(C\) to \(H\) is not fixed by the fact
that \(C\) entails \(B\); it requires a further model of how a white shoe
is observed, as Section~\ref{sec:p15-r43-rankings} records. Hempel's
objection to resolutions of this kind is set out in
Section~\ref{sec:p15-w02b-p15-024-carnap}.

\section{The Linda example in full}
\label{sec:p15-w02c-p15-027-linda}

Merin introduces the conjunction fallacy through the example he says
made it most famous, from the late 1970s
\cite[p.~273]{Merin2021Relevance}. An experimental instruction sheet
described a fictitious person, Linda, by a description \(E\) that, in
Merin's paraphrase (he says the sheet gave ``something like'' it), makes
her thirty-one, single, outspoken and bright, a philosophy graduate,
and a demonstrator for equal-opportunity causes. Participants were asked to
order by probability three claims about her: that she works as a bank
teller (\(B\)), that she is active in the feminist movement (\(F\)), and
that both hold (\(BF\)), presumably given \(E\). Merin reports that an
unexpectedly large share of respondents put \(BF\) ahead of \(B\),
against the law \(P(BF\mid E)\leq P(B\mid E)\).

His argument then has three further steps. First, the description was
constructed to fit \(F\) and not to fit \(B\); in the terminology
Merin uses, it is meant to be ``representative'' of the one and
``unrepresentative'' of the other. Second, many
probabilists would explicate the representativeness of a property
\(\Phi\) as the relevance of the statement that Linda is \(\Phi\) to the
reference description \(E\); on that explication \(F\) is positive to
\(E\) and \(B\) negative. Third, if respondents answered the request for
a probability ranking with a covert relevance ranking, the reported
ranking is feasible and natural. Merin points to Tentori, Crupi and
Russo (2013) and to Cevolani and Crupi (2015) for relevance-based
solutions.

The explication fixes a direction: the property statement is the source
and the description \(E\) the target. With the target fixed, the
ranking of sources is the task of Merin's Theorem~9, so it does not
depend on which POP-D measure is used.

\begin{hbtheorem}{thm:p15-w02c-p15-027-linda}{proved=annex:p15-w02c-p15-027-linda}{Representativeness ranking of a conjunction}
Let \(0<P(E)<1\) and \(0<P(BF)\leq P(B)<1\).
\begin{enumerate}
\item \(P(BF\mid E)\leq P(B\mid E)\).
\item For every measure satisfying POP-D,
  \(R(BF\Yright E)_P\bowtie R(B\Yright E)_P\) exactly when
  \(P(E\mid BF)\bowtie P(E\mid B)\), where \(\bowtie\) is any of
  \(<,=,>\). The same comparison is made by
  \(R_Q(BF\Yright E)_P\) and \(R_Q(B\Yright E)_P\).
\item If \(P(B\neg F)>0\), then \(P(E\mid BF)>P(E\mid B)\) exactly when
  \(P(E\mid BF)>P(E\mid B\neg F)\).
\item Take the model of Section~\ref{sec:p15-r43-use}: \(P(E)=1/10\),
  and the cells \(BF\), \(B\neg F\), \(\neg BF\), \(\neg B\neg F\) have
  prior masses \(1/50\), \(9/50\), \(12/25\), \(8/25\) and conditional
  masses given \(E\) of \(1/10\), \(1/20\), \(7/10\), \(3/20\). Then
  \[
   P(E\mid BF)=\tfrac12,\quad P(E\mid F)=\tfrac{4}{25},\quad
   P(E\mid B)=\tfrac{3}{40},
  \]
  so \(F\) is positive and \(B\) negative to \(E\), and
  \(R_D(BF\Yright E)_P=2/5\), \(R_D(F\Yright E)_P=3/50\),
  \(R_D(B\Yright E)_P=-1/40\), while
  \(P(BF\mid E)=1/10<3/20=P(B\mid E)\).
\end{enumerate}
\end{hbtheorem}

Item~2 is what makes the covert-relevance reading measure-robust in
Merin's direction. Chapter~\ref{sec:p15-rationality-4-3} computes the
same model in the opposite direction, with \(E\) as source and the
properties as targets; there the comparison is a target ranking, for
which Theorem~\ref{thm:p15-r43-rankings} shows that D, Q and L can
disagree. Item~3 says exactly what the reading requires of the
description: among bank tellers, \(E\) must be likelier for feminists
than for nonfeminists. That is a condition on the respondent's
probabilities, not on the probability law, which item~1 states holds
without exception.

\begin{hbwarning}{def:p15-w02c-p15-027-linda-standing}{What the Linda report establishes}
The empirical part of the example has the following standing in
Merin's chapter.
\begin{description}
\item[Task.] Ranking three probabilities after reading a personality
  description; the description is paraphrased, not quoted, and the
  response format and any other options are not given.
\item[Population and design.] Not stated. Merin dates the example to the
  late 1970s and gives no reference for the original study in the
  passage; the design is described only by the claim that \(E\) was
  built to be representative of \(F\) and unrepresentative of \(B\).
\item[Measure and uncertainty.] The proportion of respondents ranking
  \(BF\) above \(B\), reported only as surprisingly large. No count,
  sample size, interval or test is given.
\item[Competing explanations.] The chapter offers one, covert relevance
  judgment, and treats it as a possibility (``if one assumes''). It
  does not discuss alternatives. Accounts of the original tasks and of
  other explanations, including linguistic readings of the conjunction,
  are given in Chapters~\ref{sec:p19-conjunction-judgments}
  and~\ref{sec:p19-conjunction-language}.
\item[Replication and generalization.] Not addressed.
\end{description}
The report therefore supports no frequency, effect size or population
claim here. Theorem~\ref{thm:p15-w02c-p15-027-linda} establishes only
that the covert-relevance reading is coherent and what it requires; it
does not show that respondents performed relevance judgments, and the
reading does not make a ranking of \(BF\) above \(B\) a correct
probability judgment. The attribution of the representativeness
explication to many probabilists is a report about the literature for
which the chapter cites no specific source.
\end{hbwarning}

\section{Base-rate neglect as Merin states it}
\label{sec:p15-w02c-p15-027-base-rate}

Merin's psychological section opens with a general claim: in some
contexts, experimental subjects asked for a posterior probability or for
posterior odds appear to give replies that are consistent with their
being covert relevance judgments. He names two prominent instances, the
neglect of prior probabilities or base rates and the conjunction
fallacy, and refers in note~19 to Kahneman (2011, chapters~15--16) and to
the chapters by Evans and by Hertwig and Kozyreva in the same handbook
\cite[p.~273 and p.~274, n.~19]{Merin2021Relevance}. The second instance
is the Linda example of Section~\ref{sec:p15-w02c-p15-027-linda}.

For the first instance he states a design and an argument
\cite[p.~273]{Merin2021Relevance}. The design has three elements. The
experimenter's instruction sheet supplies information about a hypothesis
\(H\); an observation \(E\) then bears on \(H\); and the participant is
asked for the posterior probability or posterior odds of \(H\). The
argument has three steps.
\begin{enumerate}
\item The neglect phenomena, when they occur, can be seen as instances of
  Bayes' theorem in odds form, his equation~(3):
  \(O(H\mid E)=O(H)\cdot R_L(E\Yright H)_P\), where \(O\) is odds
  \cite[p.~270 and p.~274, n.~10]{Merin2021Relevance}.
\item In these instances the prior odds for \(H\) are tacitly set to one,
  that is, \(H\) and \(\neg H\) are treated as equally probable. Those
  odds do not reflect the information on the instruction sheet, which
  should warrant a skewed prior, for instance one with
  \(P(H)\gg P(\neg H)\), before the observation \(E\) arrives.
\item Arithmetically, neglecting that information makes the posterior
  odds equal to the likelihood ratio, that is, to the L value of the
  relevance of \(E\) to \(H\).
\end{enumerate}
The reply is thus read as the correct value of a different quantity: the
relevance of the observation, given in place of the requested posterior.

\begin{hbtheorem}{thm:p15-w02c-p15-027-base-rate}{proved=annex:p15-w02c-p15-027-base-rate}{Posterior odds and the likelihood ratio}
Let \(0<P(H)<1\) and \(P(E)>0\), write \(O(H)=P(H)/P(\neg H)\), and let
\(O(H\mid E)=P(H\mid E)/P(\neg H\mid E)\), with value \(+\infty\) when
\(P(\neg H\mid E)=0\).
\begin{enumerate}
\item \(O(H\mid E)=O(H)\cdot R_L(E\Yright H)_P\), where
  \(O(H)\cdot(+\infty)=+\infty\).
\item \(R_L(E\Yright H)_P=O(H\mid E)\) exactly when \(O(H)=1\), that is
  \(P(H)=1/2\), or \(R_L(E\Yright H)_P\) is \(0\) or \(+\infty\). In
  probability form, \(R_L/(1+R_L)\), read as \(1\) when \(R_L=+\infty\),
  equals \(P(H\mid E)\) under exactly the same condition.
\item If \(0<R_L(E\Yright H)_P<\infty\) and \(O(H)=k\), the reply
  \(R_L(E\Yright H)_P\) differs from the posterior odds by the factor
  \(k\): it is too small when \(P(H)>P(\neg H)\) and too large when
  \(P(H)<P(\neg H)\). It depends on \(P(E\mid H)\) and \(P(E\mid\neg H)\)
  alone and is the same for every prior.
\item Let the cells \(EH\), \(\neg EH\), \(E\neg H\), \(\neg E\neg H\) have
  masses \(9/100\), \(81/100\), \(3/100\), \(7/100\). Then
  \(P(H)=9/10\), so \(O(H)=9\); \(P(E\mid H)=1/10\) and
  \(P(E\mid\neg H)=3/10\), so \(R_L(E\Yright H)_P=1/3\) and \(E\) is
  negative to \(H\); and \(O(H\mid E)=3\), so \(P(H\mid E)=3/4\). The
  reply \(1/3\) corresponds to the probability \(1/4\). The posterior
  makes \(H\) the more probable alternative, and the reply makes
  \(\neg H\) the more probable.
\end{enumerate}
\end{hbtheorem}

Item~2 is Merin's arithmetical step with its exact range: the reply
coincides with the requested posterior only when the prior is even or
the observation is conclusive. Item~3 states the size and direction of
the discrepancy, and item~4 shows that under a skewed prior it can
reverse which alternative the reply favours. Item~3 also covers a skew
in the other direction, which Merin's wording leaves open since he
offers \(P(H)\gg P(\neg H)\) only as an instance: the same arithmetic
holds when the instruction sheet warrants \(P(H)\ll P(\neg H)\). The theorem makes the covert-relevance reading
coherent. It does not show that participants compute a likelihood
ratio, and on this reading the reply is wrong as an answer to the
question asked, though correct as a relevance value.

\begin{hbwarning}{def:p15-w02c-p15-027-base-rate-standing}{What the base-rate report establishes}
The claim about base-rate neglect has the following standing in Merin's
chapter.
\begin{description}
\item[Task.] A request for the posterior probability or posterior odds of
  a hypothesis after an observation, given an instruction sheet that
  carries information bearing on the prior. No particular problem text,
  hypothesis, observation or response format is given.
\item[Population and design.] No population is named. The design is the
  three-element structure above: prior-relevant information on the
  instruction sheet, an observation, and a request for the posterior.
  The chapter cites no study for it in the passage; note~19 points to
  Kahneman's chapters and to two chapters of the same handbook, which are
  presented in Chapters~\ref{sec:p13-rationality-1-2}
  and~\ref{sec:p17-rationality-8-5}.
\item[Measure and uncertainty.] The implied measure is agreement between
  the reply and the likelihood ratio. No proportion, count, sample size,
  interval or test is reported; the claim is qualified as holding in
  some contexts and when the phenomena arise.
\item[Competing explanations.] The chapter offers one, covert relevance
  judgment, introduced with ``appear'' and ``may be seen as''. It
  discusses no alternative. The original tasks and other explanations
  are presented in Chapters~\ref{sec:p19-rationality-2-4}
  and~\ref{sec:p19-natural-frequencies-bayesian-reasoning}.
\item[Replication and generalization.] Not addressed. The chapter does
  not say in which contexts the phenomenon arises.
\end{description}
The claim has three parts of different kinds. That replies fit covert
relevance judgments is an empirical claim reported without data here.
That the instruction-sheet information should warrant a skewed prior is
a normative premise about what that information ought to do to the
prior; it is what makes setting the odds to one a neglect, and the
chapter does not argue for it beyond stating it. That neglect equates
the posterior odds with the likelihood ratio is formal and is
Theorem~\ref{thm:p15-w02c-p15-027-base-rate}. The report supports no
frequency or effect-size claim, and it is not a recommendation about how
priors should be elicited.
\end{hbwarning}

\section{Disjoint disjunctions in the boundary cases}
\label{sec:p15-w02c-p15-027-disjunction}

Merin's Theorem~10 has a single hypothesis, \(P(AB)=0\)
\cite[p.~273]{Merin2021Relevance}. Part~(a) says that the relevance of
\(A\vee B\) to a target \(C\) is then a mixture of the relevances of the
disjuncts,
\[
 R_T(A\vee B\Yright C)_P=\lambda R_T(A\Yright C)_P
   +(1-\lambda)R_T(B\Yright C)_P,\qquad 0\leq\lambda\leq1,
 \tag{\(M_\lambda\)}
\]
for \(T=D,Q\), and for \(T=L\) when neither \(A\) nor \(B\) is extremely
relevant to \(C\). Part~(b) gives the weight: \(P(A)/(P(A)+P(B))\) for D
and Q, and \(P(A\mid\neg C)/(P(A\mid\neg C)+P(B\mid\neg C))\) for L.
Theorem~\ref{thm:p15-r43-composition} proves the identity when, in
addition, \(P(A)+P(B)>0\) and the target is nondegenerate, and item~3 of
Theorem~\ref{thm:p15-w02b-p15-024-exneg-disjuncts} treats a disjunct of
probability zero beside one of positive probability. The theorem below
covers the other cases the hypothesis allows, using the
completions of Definition~\ref{def:p15-r43-measures}, and shows what
L's proviso excludes. In it the product \(\lambda\cdot(+\infty)\) is
\(+\infty\) for \(\lambda>0\) and \(0\) for \(\lambda=0\).

\begin{hbtheorem}{thm:p15-w02c-p15-027-mixture}{proved=annex:p15-w02c-p15-027-mixture}{Merin's Theorem~10 at the boundary}
Let \(P(AB)=0\) and \(T\in\{D,Q,L\}\).
\begin{enumerate}
\item \emph{Degenerate target.} If \(P(C)\in\{0,1\}\), then
  \(R_T(E\Yright C)_P=\iota_T\) for every event \(E\), so
  \((M_\lambda)\) holds for every \(\lambda\in[0,1]\). Neither disjunct is
  extremely relevant to \(C\). When \(P(A)+P(B)>0\), the D and Q weight
  of part~(b) is defined and is one of these \(\lambda\); the L weight is
  the same number when \(P(C)=0\) and is undefined when \(P(C)=1\),
  because conditioning on \(\neg C\) is then undefined.
\item \emph{Two null disjuncts.} If \(P(A)=P(B)=0\), then
  \(R_T(A\vee B\Yright C)_P=R_T(A\Yright C)_P=R_T(B\Yright C)_P=\iota_T\)
  for every \(C\), so \((M_\lambda)\) holds for every \(\lambda\in[0,1]\),
  while every weight formula of part~(b) is \(0/0\) or, for L with
  \(P(C)=1\), undefined.
\item \emph{One null disjunct.} If \(P(B)=0<P(A)\), then
  \(R_T(A\vee B\Yright C)_P=R_T(A\Yright C)_P\) and \((M_\lambda)\) holds
  with \(\lambda=1\), the D and Q weight of part~(b).
\item \emph{L's proviso.} Let \(0<P(C)<1\) and \(P(A),P(B)>0\).
  \begin{enumerate}
  \item If neither disjunct is extremely positive to \(C\), then
    \((M_\lambda)\) holds for L with the L weight of part~(b), which is
    defined; this includes the case in which a disjunct is extremely
    negative to \(C\).
  \item If \(A\,\mathrm{expos}_P\,C\) and \(B\) is not extremely positive
    to \(C\), then \(R_L(A\Yright C)_P=+\infty\),
    \(R_L(B\Yright C)_P<R_L(A\vee B\Yright C)_P<+\infty\), and no
    \(\lambda\in[0,1]\) satisfies \((M_\lambda)\) for L.
  \item If both disjuncts are extremely positive to \(C\), all three L
    values are \(+\infty\) and \((M_\lambda)\) holds for every
    \(\lambda\in[0,1]\), while the L weight of part~(b) is \(0/0\).
  \end{enumerate}
\end{enumerate}
\end{hbtheorem}

With this theorem, part~(a) holds in every case its hypothesis admits:
unconditionally for D and Q, and for L exactly where the proviso keeps
one disjunct alone from being extremely positive; item~4(b) shows that
the proviso cannot be dropped, and item~4(a) that it is needed only
against extreme positivity. Part~(b) supplies a weight in fewer cases than
part~(a) holds: where both disjuncts are null, and for L where the target
is certain or both disjuncts are extremely positive, the identity holds
with any weight and the formula gives none. Item~3 is the mirror image,
under exchange of \(A\) and \(B\), of item~3 of
Theorem~\ref{thm:p15-w02b-p15-024-exneg-disjuncts}; for L its weight is
one whenever \(P(C)<1\) and \(P(A\mid\neg C)>0\).

\section{Conjunctive evidence in source position}
\label{sec:p15-w02c-p15-027-conjunction}

Merin's Lemma~3 says that if \(E_1\) and \(E_2\) are independent
conditionally on \(H\) and on \(\neg H\), then the L value of the
conjunction \(E_1E_2\) for the target \(H\) is the product of the two L
values, and its LL value is the sum of the two LL values
\cite[p.~271]{Merin2021Relevance}. The independence hypothesis is the
only condition stated. Theorem~\ref{thm:p15-r43-composition} proves the
identities when the likelihood ratios are positive and finite, and
Section~\ref{sec:p15-w02b-p15-024-lemma3} records the misprinted
subscript in the printed sum. The hypothesis allows further cases: a
source may exclude or entail \(H\) almost surely, so that its L value is
\(0\) or \(+\infty\), and a source or the conjunction may have
probability zero, where the completion clause of
Definition~\ref{def:p15-r43-measures} assigns the neutral value. The
next theorem settles every such case. In it, products in \([0,+\infty]\)
and sums in \([-\infty,+\infty]\) follow the rules
\(c\cdot(+\infty)=(+\infty)\cdot c=+\infty\) for \(0<c\leq+\infty\),
\(c+(+\infty)=(+\infty)+c=+\infty\) for \(c>-\infty\) and
\(c+(-\infty)=(-\infty)+c=-\infty\) for \(c<+\infty\); the product
\(0\cdot(+\infty)\) and the sum \((+\infty)+(-\infty)\) are left
undefined.

\begin{hbtheorem}{thm:p15-w02c-p15-027-lemma3}{proved=annex:p15-w02c-p15-027-lemma3}{Merin's Lemma~3 in every case its hypothesis allows}
For events \(E_1,E_2,H\) and \(i=1,2\) write
\(l_i=R_L(E_i\Yright H)_P\) and \(m_i=R_{LL}(E_i\Yright H)_P\), and write
\(l_{12}=R_L(E_1E_2\Yright H)_P\) and \(m_{12}=R_{LL}(E_1E_2\Yright H)_P\).
\begin{enumerate}
\item If \(P(H)\in\{0,1\}\), then \(l_1=l_2=l_{12}=1\) and
  \(m_1=m_2=m_{12}=0\), so \(l_{12}=l_1l_2\) and \(m_{12}=m_1+m_2\) for
  all \(E_1,E_2\). Definition~\ref{def:p15-r43-extreme} writes
  \((E_1\perp E_2\mid\pm H)_P\) only when \(0<P(H)<1\); this item shows
  that the identities hold here however the hypothesis is read.
\end{enumerate}
Let \(0<P(H)<1\) and \((E_1\perp E_2\mid\pm H)_P\) in the remaining items.
\begin{enumerate}
\setcounter{enumi}{1}
\item If \(P(E_1E_2)>0\), then \(\{l_1,l_2\}\neq\{0,+\infty\}\), the
  product \(l_1l_2\) and the sum \(m_1+m_2\) are defined, and
  \[
   l_{12}=l_1l_2,\qquad m_{12}=m_1+m_2 .
  \]
  In particular \(l_{12}=0\) when a source is extremely negative to
  \(H\), and \(l_{12}=+\infty\) when a source is extremely positive to
  \(H\) and the other is not extremely negative to it.
\item If \(P(E_1E_2)=0\), then \(l_{12}=1\) and \(m_{12}=0\), and exactly
  one of the following holds.
  \begin{enumerate}
  \item \(P(E_1)=0\) or \(P(E_2)=0\). A null source has L value \(1\)
    and LL value \(0\), and the two identities hold if and only if the
    other source is null or irrelevant to \(H\).
  \item \(P(E_1)>0\) and \(P(E_2)>0\). Then one source is extremely
    positive and the other extremely negative to \(H\),
    \(\{l_1,l_2\}=\{0,+\infty\}\) and \(\{m_1,m_2\}=\{-\infty,+\infty\}\),
    so neither \(l_1l_2\) nor \(m_1+m_2\) is defined.
  \end{enumerate}
\item The failure in item~3(a) does not depend on the completion clause.
  If \(R_T\) satisfies Quality Representation
  (Definition~\ref{def:p15-w02c-p15-027-measure}) with neutral value
  \(\iota_T\), \(P(E)=0\) and \(E'\) is any event, then
  \(R_T(E\Yright H)_P=R_T(EE'\Yright H)_P=\iota_T\). So
  \(R_T(EE'\Yright H)_P=R_T(E\Yright H)_P\cdot R_T(E'\Yright H)_P\) fails
  when \(\iota_T=1\), and
  \(R_T(EE'\Yright H)_P=R_T(E\Yright H)_P+R_T(E'\Yright H)_P\) fails when
  \(\iota_T=0\), whenever \(E'\) is relevant to \(H\).
\item \emph{Examples.} (i) Let \(P(H)=1/2\), let \(E_1\) be null, and let
  \(P(E_2\mid H)=1/2\) and \(P(E_2\mid\neg H)=1/4\). The hypothesis holds,
  \(l_{12}=1\) and \(m_{12}=0\), while \(l_1l_2=2\) and
  \(m_1+m_2=\ln2\). (ii) On four points of mass \(1/4\) each, let \(H\)
  be points \(1,2\), \(E_1\) point \(1\) and \(E_2\) point \(3\). The
  hypothesis holds, \(l_1=+\infty\), \(l_2=0\), \(m_1=+\infty\),
  \(m_2=-\infty\), and \(l_{12}=1\), \(m_{12}=0\).
\end{enumerate}
\end{hbtheorem}

By items~1 to~3, Merin's Lemma~3, read with the completions of
Definition~\ref{def:p15-r43-measures}, holds exactly in these cases: the target is
certain or impossible; the conjunction has positive probability; or one
source is null and the other is null or irrelevant to \(H\). In every
other case its hypothesis allows, either the printed identity is false
(item~3(a) with a relevant second source) or its right side has no value
(item~3(b)). The statement on p.~271 therefore needs a further
condition; the natural one is \(P(E_1E_2)>0\), under which item~2 gives
both identities with extreme values included. This corrects the
source's unrestricted statement without affecting the cases in which the
identities serve to pool observed evidence: a conjunction that is
observed and conditioned on has positive probability. Item~4 shows that no other choice of neutral completion
would rescue the null case, since any measure that represents
irrelevance by a fixed value must give a null source and a null
conjunction that same value.

Merin's paragraph on conditional compositionality for conjunction makes
three further claims \cite[p.~272]{Merin2021Relevance}. First, Lemma~3
exhibits, for L and LL, the gold standard for useful conjunctive
compositionality in source position. Second, the qualitative
Theorem~6, Reichenbach's result on screened positive links
(Theorem~\ref{thm:p15-r43-screening}), anticipates its usefulness.
Third, no other familiar measure quite matches it; for a stark contrast
he points to Carnap's covariance measure, which is additive for a
conjunction when \(P(A\vee B)=1\), a composition he calls neither
practically useful nor intuitive. The next theorem gives the third
claim an exact form for the measures of Definition~\ref{def:p15-r43-measures}.

\begin{hbtheorem}{thm:p15-w02c-p15-027-pooling}{proved=annex:p15-w02c-p15-027-pooling}{Conjunctive composition for D, Q, QL and Z needs the prior}
Let \(0<h=P(H)<1\), \((E_1\perp E_2\mid\pm H)_P\) and \(P(E_1E_2)>0\), and
write \(q_i=P(H\mid E_i)\) and \(q_{12}=P(H\mid E_1E_2)\).
\begin{enumerate}
\item The denominator below is positive, and
  \[
   q_{12}=\frac{q_1q_2(1-h)}{q_1q_2(1-h)+(1-q_1)(1-q_2)h}.
  \]
\item For \(T\in\{D,Q,QL,Z\}\), the prior \(h\) and the value
  \(r=R_T(E\Yright H)_P\) of a source with \(P(E)>0\) determine
  \(q=P(H\mid E)\): \(q=h+r\) for D, \(q=hr\) for Q, \(q=he^{r}\) for QL
  (with \(e^{-\infty}=0\)), and for Z, \(q=h+(1-h)r\) if \(r>0\) and
  \(q=h+hr\) if \(r\leq0\). Hence \(h\) and the two constituent values
  determine \(R_T(E_1E_2\Yright H)_P\).
\item The two constituent values alone do not determine it. Let
  \(M(h,a,b)\) be the model on which \(P(H)=h\), and \(E_1\), \(E_2\) are
  independent given \(H\), each with probability \(a\), and independent
  given \(\neg H\), each with probability \(b\). Then
  \[
  \begin{array}{l|ccc}
   \text{model} & T & R_T(E_i\Yright H)_P & R_T(E_1E_2\Yright H)_P\\\hline
   M(1/2,3/4,1/4) & D & 1/4 & 2/5\\
   M(1/4,3/4,1/4) & D & 1/4 & 1/2\\
   M(1/2,3/4,1/4) & Q & 3/2 & 9/5\\
   M(1/3,1/2,1/4) & Q & 3/2 & 2\\
   M(1/2,3/4,1/4) & Z & 1/2 & 4/5\\
   M(1/4,5/6,1/6) & Z & 1/2 & 6/7
  \end{array}
  \]
  and the QL values are the logarithms of the Q values. So for each
  \(T\in\{D,Q,QL,Z\}\) there is no function \(g_T\) with
  \(R_T(E_1E_2\Yright H)_P=g_T\bigl(R_T(E_1\Yright H)_P,R_T(E_2\Yright H)_P\bigr)\)
  on all models meeting the hypotheses.
\item By contrast the L and LL values of the conjunction are the product
  and the sum of the constituent values
  (Theorem~\ref{thm:p15-w02c-p15-027-lemma3}, item~2), and on models
  whose likelihoods lie strictly between zero and one the KO value is
  \((k_1+k_2)/(1+k_1k_2)\) in the constituent values \(k_1,k_2\)
  (Theorem~\ref{thm:p15-w02b-p15-024-ko-transform}). In both
  \(M(1/2,3/4,1/4)\) and \(M(1/4,3/4,1/4)\) the constituent L values are
  \(3\) and the conjunction's is \(9\), while the D values of the
  conjunction differ.
\end{enumerate}
\end{hbtheorem}

The three claims have different standing. The first is an evaluation.
What can be proved is the feature that makes the rule useful: by
item~2 of Theorem~\ref{thm:p15-w02c-p15-027-lemma3} and the odds form of
Theorem~\ref{thm:p15-r43-measures}, the posterior odds of \(H\) after two
screened reports are the prior odds times two numbers each computed
from one report alone, so the reports can be pooled before the prior is
known and the prior enters once. Item~3 of
Theorem~\ref{thm:p15-w02c-p15-027-pooling} shows that D, Q, QL and Z lack
this: whoever pools their values must also know the prior. That this
makes L and LL the standard against which others are measured is
Merin's judgment; the chapter gives no criterion of usefulness beyond
the remark that most constraints under which homomorphisms hold are of
no practical use (Theorem~\ref{thm:p15-w02c-p15-027-homomorphism}). The
second claim relates two results. Its formal content is that the
conclusion of Theorem~6 for screened positive sources follows from the
product rule, since a product of two numbers greater than one is greater
than one; this derivation is item~2 of
Theorem~\ref{thm:p15-w02b-p15-024-diagnostic-additivity}. That the
qualitative theorem \emph{presages} the quantitative one is a reading
of that relation, not a further theorem. The third claim is comparative
and Merin gives no proof of it. On the reading that no other familiar
measure composes conjunctive evidence from the constituent values alone
by multiplication or addition, it is true:
Theorem~\ref{thm:p15-w02c-p15-027-pooling} proves it for D, Q, QL and Z,
and Theorem~\ref{thm:p15-w02b-p15-024-ko-transform} for KO, which has a
rule in the constituent values alone that is neither a sum nor a
product. Merin's \emph{quite} fits KO: it comes close, since it needs no
prior, but its rule is the addition law of the hyperbolic tangent. The
Carnap contrast is proved in Theorem~\ref{thm:p15-r43-composition}.

\section{The chapter's claims by kind}
\label{sec:p15-w02c-p15-027-claims}

Convention~\ref{def:p15-w02b-p15-024-claim-kinds} sorts Merin's claims
into formal, normative or methodological, historical and empirical
ones, and names three empirical claims. The chapter contains more of
each non-formal kind than that list records. Its normative and
justificatory claims, with the formal results that bear on them, are
these.
\begin{enumerate}
\item \emph{Scope and vagueness of the informal notion} (p.~265). The
  same relations of speaking for and against hold between statements
  about past events and between possible experimental results and
  timeless scientific hypotheses; in legal and scientific usage, to
  speak for \(B\) is to be evidence for \(B\); and the
  informal idea, plausible as it is, stays vague and prone to apparent
  exceptions, so it needs a framework. This is the motivation for the
  explication of Definition~\ref{def:p15-r43-polarity}; no argument
  beyond the weather example is given for it.
\item \emph{Conclusive relevance and conditionals} (p.~266). Extreme
  positivity explicates what an unreserved assertion of \emph{if \(A\),
  then \(B\)} conveys, and what everyday speech mostly means by
  entailment; Merin qualifies the first point by noting that the
  conditional can also be asserted with less certainty. The proposal and
  its limit are set out in Section~\ref{sec:p15-w02c-p15-027-extreme},
  with the independence of the two conditions in
  Theorem~\ref{thm:p15-w02c-p15-027-transitivity}; the proposal is an
  explication, not supported by data about speakers.
\item \emph{Notation and terminology} (p.~268). A directed,
  content-neutral notation is recommended, and the word \emph{relevance}
  is preferred to \emph{confirmation}, \emph{corroboration} and
  \emph{support} on two grounds: those words need the prefix
  \emph{incremental} to avoid being read as \(P(H\mid E)\), and they sound
  awkward with the prefix \emph{negative}, whereas \emph{relevance}, the
  word of Keynes (1921), Carnap (1950) and Jeffrey (2004), has neither
  problem. The four-step argument is given in
  Section~\ref{sec:p15-w02c-p15-027-direction}; only its premise about
  direction has a formal basis
  (Theorem~\ref{thm:p15-w02c-p15-027-direction}), and the grounds about
  misreading and awkwardness are claims about usage. Two remarks go with
  it: that a content-neutral directed notation lets one investigate
  freely what motivates the conventional choice of direction (p.~268),
  a methodological claim whose investigation the chapter does not carry
  out in general; and that Merin mostly keeps the format
  \(R(E\Yright H)\) for cross-reference to the confirmation literature,
  as \(A\) and \(B\) serve the content-neutral mathematical literature
  (p.~269), a choice of presentation. Both are discussed after the
  four steps in Section~\ref{sec:p15-w02c-p15-027-direction}.
\item \emph{What ordinal equivalence leaves open} (p.~270). Distinct but
  ordinally equivalent measures may differ in the convenience of their
  ranges and in their algebraic properties, the latter being the more
  important for cognitive-engineering applications. The differences are
  exhibited in Section~\ref{sec:p15-w02c-p15-027-ordinal}; the ranking of
  the two is Merin's judgment, given without an application.
\item \emph{Amount symmetry} (p.~272). Intuitions on amount symmetry
  are strong only for asymmetric extreme positivity. In the
  divisibility example, D separates the two directions while Q does
  not, and Merin predicts that many will balk at the Q verdict because
  it ignores the asymmetry of entailment. The formal facts are
  Theorem~\ref{thm:p15-r43-measures} and
  Theorem~\ref{thm:p15-w02c-p15-027-direction}; the claims about which
  intuitions are strong and how readers will react are reported without
  data.
\item \emph{Extreme valuation} (p.~272). L, Z and their ordinal
  equivalents take the extreme values of their ranges exactly at
  extreme relevance, and Merin says this pleases anyone for whom
  conclusive argument is the strongest argument. The property is stated
  in Section~\ref{sec:p15-r43-measures}; the evaluation is conditional
  on that preference and is not argued further.
\item \emph{Homomorphisms} (p.~272). Structure-preserving maps from
  Boolean operations on propositions to simple arithmetical operations on
  relevance values (inversion, addition, weighted mean) are claimed to be
  essential if the relevance of compound propositions is to be
  predictable: for many of their ordinal properties and for all of their
  quantitative ones. Every relevance measure
  has some such homomorphism under narrow enough constraints; most of
  these constraints are of no practical use, and not every homomorphism
  is desirable. The three instances Merin discusses are proved in
  Theorems~\ref{thm:p15-r43-measures}
  and~\ref{thm:p15-w02c-p15-027-negation} (negation) and
  Theorem~\ref{thm:p15-r43-composition} (conjunction under conditional
  independence, with every case of the hypothesis in
  Theorem~\ref{thm:p15-w02c-p15-027-lemma3}; disjunction of disjoint
  sources, with the boundary cases in
  Theorem~\ref{thm:p15-w02c-p15-027-mixture}); Carnap's additivity
  for conjunction when \(P(A\vee B)=1\), in the last theorem, is his
  example of one that is neither useful nor intuitive. The general
  existence claim is Theorem~\ref{thm:p15-w02c-p15-027-homomorphism}
  below.
\item \emph{Seeking independencies} (p.~273). Because the qualitative
  transitivity results and the composition laws depend on conditional
  independence, Merin advises looking for conditional independencies
  in whatever domain the theory is applied to. The formal results
  (Theorems~\ref{thm:p15-r43-screening} and~\ref{thm:p15-r43-composition})
  show what independence yields; whether it holds in a domain is an
  empirical question that the recommendation does not settle.
\item \emph{Choice of measure} (p.~274). LL as the default, Q and QL
  when direction is unspecified, and pluralism for special purposes, as
  recorded at the end of Chapter~\ref{sec:p15-rationality-4-3}.
\item \emph{The product form and the symmetry gloss} (p.~266).
  Definition~3 is preferred to Definition~2 because it needs no
  completion clause and shows the symmetry of the relations; and
  Theorem~1 is glossed intuitively as saying that if \(A\) speaks for
  \(B\), then \(B\) speaks for \(A\), however weakly. Both reasons are
  proved in Theorem~\ref{thm:p15-w02c-p15-027-product}, which also gives
  the gloss a numerical instance; the gloss itself is a reading, not a
  theorem.
\item \emph{Unconditional relevance as a special case} (p.~267).
  Definition~3 is presented as the case of Definition~5 with a
  tautological condition, on the ground that \(P(A)=P(A\mid\top)\).
  This is item~5 of Theorem~\ref{thm:p15-w02c-p15-027-product}.
\item \emph{What the characterization determines} (p.~271). The
  \emph{if} parts of Theorem~7 yield classes of ordinally equivalent
  measures, and to determine the specific measures further conditions
  are needed, smoothness among them. Theorem~\ref{thm:p15-w02c-p15-027-classes} shows
  that smoothness excludes some members of the class of D but does not
  by itself determine D; the chapter names no other condition.
\item \emph{The skewed prior} (p.~273). Information on an experimenter's
  instruction sheet should warrant a skewed prior before the observation.
  This premise is what makes the even prior odds of
  Section~\ref{sec:p15-w02c-p15-027-base-rate} a neglect; it is stated,
  not argued.
\item \emph{Conditional compositionality for conjunction} (p.~272). Lemma~3
  is called the gold standard for useful conjunctive composition in
  source position, Theorem~6 is said to presage its usefulness, and no
  other familiar measure is said to quite match it. The first is an
  evaluation, the second a reading of how two results are related, and
  the third a comparative formal claim that Merin asserts without proof.
  Their standing, and the proof of the third in an exact form, are given
  in Section~\ref{sec:p15-w02c-p15-027-conjunction}
  (Theorems~\ref{thm:p15-w02c-p15-027-lemma3}
  and~\ref{thm:p15-w02c-p15-027-pooling}).
\item \emph{The logarithmic transforms} (p.~269). QL and LL are called
  the most useful transforms of Q and L. No reason is given at that
  point; the equations of the next page (p.~270) show what the
  logarithms do: they turn the multiplicative updates
  \(q=hR_Q\) and \(O(H\mid E)=O(H)R_L\) into additive ones
  (Theorem~\ref{thm:p15-r43-measures}), and LL turns the product of
  Lemma~3 into a sum (Theorem~\ref{thm:p15-w02c-p15-027-lemma3}). Since
  each is an increasing transform of its parent, it changes no ranking
  (Theorem~\ref{thm:p15-w02c-p15-027-ordinal}); what it changes is the
  arithmetic, the kind of difference item~4 above calls the more
  important. QL has no prior-free rule for screened conjunctions
  (Theorem~\ref{thm:p15-w02c-p15-027-pooling}), so its usefulness rests on
  the update form alone. \emph{Most useful} is Merin's evaluation; apart
  from the remark on KO (p.~271), the chapter compares the logarithms
  with no other transform of Q or L.
\end{enumerate}

\begin{hbtheorem}{thm:p15-w02c-p15-027-homomorphism}{proved=annex:p15-w02c-p15-027-homomorphism}{Every relevance measure has a conditional conjunction homomorphism}
Let \(R_T\) satisfy Definition~\ref{def:p15-w02c-p15-027-measure} with
neutral value \(\iota_T\). If \(P(B)=1\), then for every source \(A\)
and target \(H\), \(R_T(B\Yright H)_P=\iota_T\) and
\(R_T(AB\Yright H)_P=R_T(A\Yright H)_P\). Hence
\(R_T(AB\Yright H)_P=R_T(A\Yright H)_P+R_T(B\Yright H)_P\) when
\(\iota_T=0\), and
\(R_T(AB\Yright H)_P=R_T(A\Yright H)_P\cdot R_T(B\Yright H)_P\) when
\(\iota_T=1\) and \(R_T(A\Yright H)_P\) is finite.
\end{hbtheorem}

The constraint, that one conjunct is certain, is exactly the kind Merin
calls practically useless: a certain conjunct carries no evidence. The
theorem shows that the bare existence of a conditional homomorphism says
nothing about a measure's merits, which is why his assessment rests on
the three instances whose constraints occur in use.

The empirical and descriptive claims of the chapter, each reported
without a study, population, design or estimate of uncertainty unless
stated, are as follows. Claims about people:
\begin{enumerate}
\item Asked for a posterior probability or posterior odds, experimental
  subjects sometimes give answers that fit a hidden relevance judgment
  instead, the leading cases being base-rate neglect and the conjunction
  fallacy (p.~273). For base-rate neglect Merin states the requested
  quantity, the neglected instruction-sheet information and the
  arithmetic, with the standing given in
  Warning~\ref{def:p15-w02c-p15-027-base-rate-standing}; for the
  conjunction fallacy, the Linda task of item~2.
\item In the Linda task an unexpectedly large share of respondents put
  \(BF\) ahead of \(B\) (p.~273), with the standing given in
  Warning~\ref{def:p15-w02c-p15-027-linda-standing}.
\item There is evidence that \(\delta p\) captures how participants
  judge the bearing of \(E\) on \(H\), reported by
  Skovgaard-Olsen, Singmann and Klauer (2016b); their 2016a paper
  investigated its relation to indicative conditionals (p.~273). No
  task, sample or result is given.
\item Relevance theory conjectures that maximizing relevance, in the
  sense of greater cognitive effects at lower processing cost, is a
  general principle of cognition and communication; it has been applied
  to the Wason selection task by Sperber, Cara and Girotto (1995)
  (p.~273). Merin records that this notion of relevance differs from
  evidential relevance.
\item Many will balk at the symmetry of Q in the divisibility example
  (p.~272); intuitions on amount symmetry are strong only for
  asymmetric extreme positivity (p.~272).
\item On the common view, a hearer draws the defeasible not-both
  implicature of \emph{A or B} from the speaker's failure to assert
  \emph{A and B} (p.~274, n.~16).
\end{enumerate}
Claims about the literature and about practice:
\begin{enumerate}
\item[7.] The literature commonly describes a degree of relevance as an
  amount of incremental confirmation, written \(C_P(H,E)\) or
  \(C_P(E,H)\) (p.~268); likelihood glosses are often counterfactual
  (p.~274, n.~5).
\item[8.] All measures discussed in the literature satisfy
  Definition~8, and measures in use have neutral value zero or one
  (p.~269).
\item[9.] Of the relevance measures, D, Q and L have the longest
  history and the widest use, and QL and LL are the most common
  transforms of Q and L (p.~269); statisticians use LL widely
  (p.~270).
\item[10.] POP-D seems mandatory to most commentators, and LIK-D is
  often called the law of likelihood (p.~271); philosophers, unlike
  statisticians, extend the likelihood concept to hypotheses that are
  compatible and not exhaustive (p.~274, n.~14).
\item[11.] Professed opinion usually commends inversion under negation
  in target position and rejects it in source position, citing Eells
  and Fitelson (2002) (p.~272).
\item[12.] Many probabilists would explicate representativeness as
  relevance to a reference description (p.~273).
\item[13.] Relevance theory is popular in linguistics (p.~273).
\item[14.] Statisticians facing known data typically rely on existing
  models of conditional frequencies (p.~268), and some philosophers
  recommend measure pluralism (p.~274).
\item[15.] Pearl's thesis that diagnostic categories are formed to
  screen off symptoms (p.~267), and the general-relativity case, where
  the entailment of Mercury's orbit was widely judged to raise
  confidence (p.~268).
\item[16.] The most familiar applications of the symmetry of relevance
  and of its reversal under negation presuppose a sorting of
  propositions into hypotheses or theories and data reports or evidence
  statements (p.~267), a sorting that Merin then calls often dependent on
  the application (p.~268); see
  Section~\ref{sec:p15-w02c-p15-027-extreme}.
\end{enumerate}
Items~3 and~4 name their sources; the rest are Merin's own reports of
opinion, usage or practice. Items~9 and~11 concern the literature
Merin cites; for D, Q and L the historical attributions are compared in
Section~\ref{sec:p15-w02b-p15-024-history}. None of the sixteen is
used as a premise of a formal result here.

The historical claims are of two kinds. The first are attributions to
earlier authors: those compared in
Section~\ref{sec:p15-w02b-p15-024-history}, the date of the Linda example
above, the 2016-election instance of the transitivity counterexample
(Section~\ref{sec:p15-w02c-p15-027-extreme}), and Steel's argument that
POP-D is necessary for enumerative induction (p.~271), listed among the
arguments attributed but not given in
Proofboundary~\ref{def:p15-w02b-p15-024-dependents}. The second are
Merin's claims about what his own chapter adds.
\begin{enumerate}
\item In its summary the chapter says that, on the ordinal and
  quantitative side of relevance, it offers a notational proposal and
  two extensions of the state of the art: conditions, relative to a
  task, under which all relevance measures meeting an intuitive
  desideratum that the major measures share are ordinally equivalent, which gives many comparative
  judgments a useful degree of independence from the measure; and
  relevance properties of disjunctive evidence, to complement known facts
  about conjunctive evidence (p.~265). The notational proposal is the
  signpost notation of his Definition~6 (p.~268); the desideratum is
  POP-D; the first extension is his Theorems~8 and~9 and Lemma~4, and the
  second his Theorem~10.
\item Introducing Theorems~8 and~9 and Lemma~4, Merin presents them as
  his own results, offered to locate where the known pairwise ordinal
  inequivalence of D, Q and L shows itself (p.~271). The conclusion
  repeats both extensions: mutual extreme negativity of the disjuncts is
  described as a condition for the compositionality of disjunctive
  evidence that was established for the major measures, and the
  task-relative equivalence result as an assurance that many ordinal
  judgments are not unreasonably measure dependent (p.~274).
\item The summary describes the chapter in two further ways (p.~265). It
  says that the abstract theory of qualitative relevance is explored
  quite fully, notably the role of conditional irrelevance; and it says
  that the central role of relevance for reasoning is emphasized
  throughout. The qualitative theory in the chapter consists of
  Definitions~1 to~5, Theorems~1 to~6 and Lemmas~1 and~2, with the
  application of section~4 (pp.~266--268). Conditional irrelevance
  carries Theorems~5 and~6, Lemma~2, the drug-and-cure example and the
  test of Pearl's thesis (p.~267). Most of these results are stated
  without proof: the proofs of Theorems~1 and~3 are indicated in a
  sentence each (p.~266), and Theorems~5 and~6 are attributed to
  Reichenbach (1956) (p.~267 and p.~274, n.~3). \emph{Quite fully} is
  Merin's own assessment; the chapter gives no list of qualitative
  properties against which it could be checked, and the proofs, models
  and counterexamples are supplied in
  Chapter~\ref{sec:p15-rationality-4-3} and in
  Sections~\ref{sec:p15-w02c-p15-027-product}
  and~\ref{sec:p15-w02c-p15-027-extreme}. The emphasis on reasoning is a
  description of the chapter's own framing, borne out by its opening on
  speaking for and against in argument (p.~265), its section~4 on
  qualitative relevance reasoning (pp.~267--268), its reading of
  relevance measures as quasi-utilities serving the rhetorical end of
  establishing a target (p.~273), its psychological section (p.~273) and
  its conclusion (p.~274); it is not a claim that admits a formal test.
\end{enumerate}
The chapter cites no earlier source for Theorems~8 to~10 or Lemma~4;
note~16 cites Merin (1999) for one ordinal result that Theorem~9
generalizes. The priority claims are reported as Merin's; this book
does not assess them against the earlier literature. The mathematical
content is proved independently: Theorems~8 and~9 in
Theorem~\ref{thm:p15-r43-rankings}, Lemma~4 in
Theorem~\ref{thm:p15-w02b-p15-024-direction}, Theorem~10 in
Theorem~\ref{thm:p15-r43-composition} with
Theorems~\ref{thm:p15-w02b-p15-024-exneg-disjuncts}
and~\ref{thm:p15-w02c-p15-027-mixture}.

Two editorial facts remain. The chapter carries an
editorial note: Arthur Merin died on 24 May 2019, shortly after
finishing it, and the editors made the final minor revisions
\cite[p.~274]{Merin2021Relevance}. Note~15, on the white-shoe step of
the raven problem, is given in Section~\ref{sec:p15-w02c-p15-027-raven}.

\section{Relationships and sources}
\label{sec:p15-w02c-p15-027-relationships}

Theorem~\ref{thm:p15-w02c-p15-027-ordinal} is the bridge between
Merin's Definition~9 and the ordinal characterizations of
Theorem~\ref{thm:p15-r43-characterization} and
Theorem~\ref{thm:p15-w02b-p15-024-z-characterization}, which are stated
as \(R=f\circ R_T\) for an increasing \(f\): by item~2 those families
are exactly the ordinal-equivalence classes of D, Q, L and Z among
transparent measures. Theorem~\ref{thm:p15-w02c-p15-027-direction}
supplements the direction result for source rankings,
Theorem~\ref{thm:p15-w02b-p15-024-direction}, which shows that for one
target the direction does not matter under POP-D; the present theorem
shows that for one pair it can. The Linda analysis connects the
relevance mathematics with the conjunction tasks presented in
Chapters~\ref{sec:p19-conjunction-judgments}
and~\ref{sec:p19-conjunction-language}, and the raven analysis with
Carnap's and Hempel's treatments in
Section~\ref{sec:p15-w02b-p15-024-carnap}.
Theorem~\ref{thm:p15-w02c-p15-027-transitivity} states in general form
the counterexample that Theorem~\ref{thm:p15-r43-polarity} realizes in one
model, and Theorem~\ref{thm:p15-w02c-p15-027-simpson} gives a population
reading of the screening formula of
Theorem~\ref{thm:p15-r43-screening}; causal readings of such populations
need a structural causal model in the sense of
Definition~\ref{def:p16-finite-scm}, which the probability model alone
does not supply.
Theorem~\ref{thm:p15-w02c-p15-027-mixture} completes the
disjunction rule of Theorem~\ref{thm:p15-r43-composition} and
Theorem~\ref{thm:p15-w02b-p15-024-exneg-disjuncts};
Theorem~\ref{thm:p15-w02c-p15-027-lemma3} completes its conjunction rule,
and Theorem~\ref{thm:p15-w02c-p15-027-pooling} sets that rule beside
D, Q, QL and Z, as Theorem~\ref{thm:p15-w02b-p15-024-ko-transform} sets it
beside KO.
Theorem~\ref{thm:p15-w02c-p15-027-base-rate} states the arithmetic
of Section~\ref{sec:p15-r43-use} with its exact range and the premise
Merin gives for it, and Section~\ref{sec:p15-w02b-p15-024-surveys}
records the pointers of his note~19.
Theorem~\ref{thm:p15-w02c-p15-027-classes} explains why the forms
\(f_D\), \(f_Q\) and \(f_L\) of
Theorem~\ref{thm:p15-r43-characterization} carry an arbitrary increasing
\(f\), and Theorem~\ref{thm:p15-w02c-p15-027-covariance} complements
Theorem~\ref{thm:p15-w02b-p15-024-carnap-bridge} on Carnap's measure.

The source is Merin's chapter~4.3 of \emph{The Handbook of Rationality}
\cite[pp.~265--275]{Merin2021Relevance}: the summary (p.~265),
Definitions~1 to~4, Theorem~1 with its gloss, and the interpretation of
Definition~4 (p.~266), Definition~5 with the tautology remark, the
transitivity counterexample, the drug-and-cure example and the opening of
section~4 (p.~267), note~3 (p.~274), Definitions~7 and~8 and note~6
(p.~269 and p.~274), the remarks on the format \(R(E\Yright H)\) and on
QL and LL (p.~269), Lemma~3 (p.~271), the paragraph on conditional
compositionality for conjunction (p.~272), Definition~9, note~9 and the remark following
Definition~9 (p.~270 and p.~274), the
direction claim and the subsection on notation and terminology
(p.~268), Theorem~7 and the remark on its equivalence classes,
Theorems~8 and~9 and Lemma~4 (p.~271), Theorem~10 (p.~273),
the raven problem and note~15 (pp.~271--272
and p.~274), the features-in-use paragraphs, including Carnap's measure
(p.~272), the psychological applications, including base-rate neglect,
with notes~10 and~19 (p.~273 and p.~274), the conclusion (p.~274), and
the editorial note (p.~274).
The models, counterexamples and proofs of this chapter are supplied
here.

\index{evidential relevance!ordinal equivalence}
\index{evidential relevance!direction}\index{raven paradox!equivalence condition}
\index{conjunction fallacy!representativeness}
\index{evidential relevance!extreme positivity}\index{Simpson's paradox!stratified population}
\index{evidential relevance!disjunctive evidence}\index{evidential relevance!conjunctive evidence}
\index{evidential relevance!conditional irrelevance}\index{base-rate neglect!relevance reading}
\index{evidential relevance!amount symmetry}
